\documentclass[11pt,a4paper]{article}

\usepackage[a4paper,landscape,margin=0.44in]{geometry}
\usepackage[T1]{fontenc}
\usepackage{lmodern}
\usepackage{microtype}
\usepackage{amsmath}
\usepackage{booktabs}
\usepackage{tabularx}
\usepackage{longtable}
\usepackage{array}
\usepackage{enumitem}
\usepackage{float}
\usepackage{xcolor}
\usepackage{hyperref}
\usepackage{xurl}
\usepackage{seqsplit}

\definecolor{LineGray}{RGB}{95,101,107}
\definecolor{LightBlue}{RGB}{233,240,248}
\definecolor{LightYellow}{RGB}{250,244,225}
\definecolor{LightRed}{RGB}{249,232,230}
\definecolor{LinkBlue}{RGB}{35,82,124}

\hypersetup{
  colorlinks=true,
  linkcolor=LinkBlue,
  urlcolor=LinkBlue,
  pdftitle={Clean Synthetic Core Medicine Search Benchmark}
}

\newcolumntype{Y}{>{\raggedright\arraybackslash}X}
\newcolumntype{L}[1]{>{\raggedright\arraybackslash}p{#1}}
\newcommand{\code}[1]{\texttt{\seqsplit{#1}}}
\newcommand{\pp}{\,pp}

\setlength{\parindent}{0pt}
\setlength{\parskip}{0.30em}
\setlength{\tabcolsep}{4.2pt}
\renewcommand{\arraystretch}{1.10}
\setlist[itemize]{leftmargin=1.35em,itemsep=0.15em,topsep=0.20em}
\clubpenalty=10000
\widowpenalty=10000
\displaywidowpenalty=10000

\newcommand{\DatasetCases}{66,257}
\newcommand{\TargetFamilies}{16,845}
\newcommand{\RepresentedSourceRows}{66,716}
\newcommand{\RemovedSourceRows}{48,284}
\newcommand{\CollapsedOccurrences}{459}
\newcommand{\ErrorTypeLabels}{4,190}
\newcommand{\DangerCases}{158}
\newcommand{\GeneratorExtremeCases}{17,913}
\newcommand{\AFiveHOne}{98.2055\%}
\newcommand{\AFiveHTwenty}{100.0000\%}
\newcommand{\BestHOneSystem}{Algorithm 5, evidence-guided rescue}
\newcommand{\BestHOneScore}{98.2055\%}
\newcommand{\BestHTwentySystem}{Algorithm 5, evidence-guided rescue}
\newcommand{\BestHTwentyScore}{100.0000\%}
\newcommand{\AFiveVsLevenHOne}{+4.1565\pp}
\newcommand{\AFiveVsLevenHTwenty}{+0.5584\pp}
\newcommand{\ExtremeCases}{123}
\newcommand{\AFiveExtremeHOne}{68.29\%}
\newcommand{\AFiveExtremeHTwenty}{100.00\%}
\newcommand{\AFiveZeroBigramHOne}{57.02\%}
\newcommand{\AFiveWorstCategory}{Combined Errors Effective 3}
\newcommand{\AFiveWorstCategoryMisses}{576}
\newcommand{\AFiveCorrectAtOne}{65,068}
\newcommand{\AFiveHOneFailures}{1,189}
\newcommand{\AFiveHOneFailureRate}{1.7945\%}
\newcommand{\AFiveRankingFailures}{1,189}
\newcommand{\AFiveRetrievalFailures}{0}
\newcommand{\AFiveNoResults}{0}
\newcommand{\UniqueOSACases}{61,228}
\newcommand{\UniqueOSAHOne}{99.9771\%}
\newcommand{\UniqueOSAFailures}{14}
\newcommand{\NonUniqueOSAFailures}{1,175}
\newcommand{\FailureDistanceLabelDifferences}{26}
\newcommand{\UniqueScoreInversions}{14}
\newcommand{\UniqueCorrectionOverrides}{0}
\newcommand{\UniqueRankingFailures}{14}
\newcommand{\UniqueRetrievalFailures}{0}
\newcommand{\AllSystemMissCount}{0}
\newcommand{\FallbackAnyChanged}{161}
\newcommand{\FallbackAnyFixed}{14}
\newcommand{\FallbackAnyBroken}{125}
\newcommand{\FallbackAnyNet}{-111}
\newcommand{\FallbackNoResultFixed}{0}
\newcommand{\FallbackNoResultBroken}{0}
\newcommand{\ManualReviewCount}{5}
\newcommand{\FailureFamilies}{1,070}
\newcommand{\TopTenFailureShare}{2.27\%}
\newcommand{\FamilyMacroHOne}{97.5120\%}
\newcommand{\FamilyMacroHTwenty}{100.0000\%}

\newcommand{\DatasetAuditRows}{%
Original generated rows & 115,000 & Complete 34-category source benchmark. \\
Source rows represented in clean core & 66,716 & Rows retained before duplicate pair collapse. \\
Source rows outside clean core & 48,284 & Behavioral, ambiguous, exact, no-match, or unsupported rows not used in this positive-retrieval denominator. \\
Primary clean unique pairs & 66,257 & Every row receives one vote in the headline score. \\
Collapsed duplicate occurrences & 459 & Repeated query-target occurrences removed from weighting. \\
Unique corrupted queries & 66,257 & No compact query is repeated. \\
Verified target families & 16,845 & All targets resolve to the 25,066-row runtime catalog. \\
Clean category labels & 22 & Membership labels; rare rows can carry more than one. \\
Clean error-type labels & 4,190 & Exact mutation labels after corrections. \\
Rows with a cleaning note & 17,230 & Rows relabelled or retained after a documented safety comparison. \\
Dangerous gold-closer pairs & 158 & Gold target is closer than the recorded alternative; both remain visible for safety analysis. \\
}
\newcommand{\SourceTransitionRows}{%
1. Single Letter Visual Confusion & 8,000 & 7,841 & Core corruption retained \\
2. Single Letter Phonetic Confusion & 6,000 & 5,983 & Core corruption retained \\
3. Multi Char Phonetic Confusion & 5,000 & 4,991 & Core corruption retained \\
4. Ligature Confusion & 4,000 & 3,992 & Core corruption retained \\
5. Single Char Deletion Position Weighted & 5,000 & 4,966 & Core corruption retained \\
6. Single Char Insertion & 3,000 & 2,988 & Core corruption retained \\
7. Transposition Position Weighted & 3,000 & 3,000 & Core corruption retained \\
8. Truncation Doctor Abbreviation & 6,000 & 0 & Absent from positive-retrieval core \\
9. Two Error Combinations & 8,000 & 7,926 & Core corruption retained \\
10. Three Error Combinations & 5,000 & 4,958 & Core corruption retained \\
11. Speed Typing Errors & 3,000 & 2,981 & Core corruption retained \\
12. Wrong Vowels In Consonant Frame & 5,000 & 4,870 & Core corruption retained \\
13. Punctuation Whitespace Copy Paste Artifacts & 1,000 & 0 & Absent from positive-retrieval core \\
14. Case Sensitivity & 500 & 0 & Absent from positive-retrieval core \\
15. Number Word Confusion & 500 & 0 & Absent from positive-retrieval core \\
16. Embedded Form Strength Parsing & 4,000 & 0 & Absent from positive-retrieval core \\
17. Dangerous ED1 Pairs & 5,000 & 158 & Core corruption retained \\
18. Substring Traps & 3,000 & 0 & Absent from positive-retrieval core \\
19. Negative No Match Expected & 3,000 & 0 & Absent from positive-retrieval core \\
20. Contradictory Form Route & 2,000 & 0 & Absent from positive-retrieval core \\
21. Cancelled Na Drug Lookup & 1,000 & 0 & Absent from positive-retrieval core \\
22. Score Gap Ambiguity Detection & 4,000 & 0 & Absent from positive-retrieval core \\
23. OCR Letter Digit Confusion & 4,000 & 3,998 & Core corruption retained \\
24. OCR + Other Error Combined & 3,000 & 2,999 & Core corruption retained \\
25. Double Letter Reduction Expansion & 3,000 & 2,987 & Core corruption retained \\
26. Keyboard Adjacent Sampled & 2,000 & 1,924 & Core corruption retained \\
27. Four + Error Combinations & 4,000 & 7 & Small stress subset retained \\
28. Consonant Frame Wrong Vowels Heavy & 4,000 & 68 & Small stress subset retained \\
29. Multi Word Name Fragmentation & 2,000 & 0 & Absent from positive-retrieval core \\
30. Autocorrect Artifacts & 2,000 & 0 & Absent from positive-retrieval core \\
31. Exact Match Baseline & 2,000 & 0 & Absent from positive-retrieval core \\
32. Exact Match With Strength & 2,000 & 0 & Absent from positive-retrieval core \\
33. Keyboard Shift Whole Word & 500 & 0 & Absent from positive-retrieval core \\
34. Prefix Ambiguity Awareness & 1,500 & 0 & Absent from positive-retrieval core \\
}

\newcommand{\FairTransitionRows}{%
Algorithm 1, current app & 78.6841\% & 83.8583\% & +5.1743\pp & 90.6578\% & 92.2333\% & +1.5755\pp \\
Algorithm 2, external fast & 82.9251\% & 92.4053\% & +9.4803\pp & 93.9740\% & 96.8788\% & +2.9048\pp \\
Algorithm 3, rank fusion & 84.7373\% & 91.8363\% & +7.0990\pp & 95.3389\% & 97.1203\% & +1.7814\pp \\
Algorithm 4, family rescue & 85.8694\% & 97.1852\% & +11.3158\pp & 95.6490\% & 99.5155\% & +3.8666\pp \\
}

\newcommand{\CleaningNoteRows}{%
No relabelling needed & 49,027 & 74.00\% & Original label agreed with the observed final mutation. \\
combined-operation case relabelled by final effective Damerau distance & 15,851 & 23.92\% & The final string, rather than the planned chain length, sets the clean category. \\
declared subtype corrected from 3 to 2 actual vowel changes & 540 & 0.82\% & Two vowels changed even though the generator subtype said three. \\
declared subtype corrected from 3 to 1 actual vowel changes & 371 & 0.56\% & One vowel changed even though the generator subtype said three. \\
declared subtype corrected from 2 to 1 actual vowel changes & 296 & 0.45\% & One vowel changed even though the generator subtype said two. \\
gold distance 1; closest competitor distance 2 & 119 & 0.18\% & The dangerous pair stays scoreable because the gold family is strictly closer. \\
combined-operation case relabelled by final effective Damerau distance; gold distance 1; closest competitor distance 2 & 16 & 0.02\% & Several documented corrections apply to the same row. \\
declared subtype corrected from 2 to 1 actual vowel changes; declared subtype corrected from 3 to 1 actual vowel changes & 14 & 0.02\% & Several documented corrections apply to the same row. \\
}

\newcommand{\CompositionRows}{%
Effective Levenshtein & 1 & 38,664 & 58.35\% & Ordinary insertion, deletion, or replacement count. \\
Effective Levenshtein & 2 & 20,248 & 30.56\% & Ordinary insertion, deletion, or replacement count. \\
Effective Levenshtein & 3 & 6,611 & 9.98\% & Ordinary insertion, deletion, or replacement count. \\
Effective Levenshtein & 4 & 421 & 0.64\% & Ordinary insertion, deletion, or replacement count. \\
Effective Levenshtein & 5 & 313 & 0.47\% & Ordinary insertion, deletion, or replacement count. \\
Effective Damerau & 1 & 41,665 & 62.88\% & Adjacent transposition counts as one edit. \\
Effective Damerau & 2 & 18,335 & 27.67\% & Adjacent transposition counts as one edit. \\
Effective Damerau & 3 & 5,694 & 8.59\% & Adjacent transposition counts as one edit. \\
Effective Damerau & 4 & 563 & 0.85\% & Adjacent transposition counts as one edit. \\
Difficulty & Easy & 586 & 0.88\% & Generator-assigned challenge level. \\
Difficulty & Extreme & 17,913 & 27.04\% & Generator-assigned challenge level. \\
Difficulty & Hard & 33,906 & 51.17\% & Generator-assigned challenge level. \\
Difficulty & Medium & 13,852 & 20.91\% & Generator-assigned challenge level. \\
Danger & Dangerous & 158 & 0.24\% & Safety label for a close named competitor. \\
Danger & Safe & 66,099 & 99.76\% & Safety label for a close named competitor. \\
Shared bigrams & 0 Shared Bigrams & 349 & 0.53\% & Adjacent two-character anchors retained in the target. \\
Shared bigrams & 1 Shared Bigram & 1,144 & 1.73\% & Adjacent two-character anchors retained in the target. \\
Shared bigrams & 2 3 Shared Bigrams & 6,741 & 10.17\% & Adjacent two-character anchors retained in the target. \\
Shared bigrams & 4 + Shared Bigrams & 58,023 & 87.57\% & Adjacent two-character anchors retained in the target. \\
}

\newcommand{\CategoryRows}{%
Single Letter Visual Confusion & 7,841 & 11.83\% & \code{happinessbarylife} & \code{HAPPINESS BABY LIFE} & 1 \\
Combined Errors Effective 2 & 7,463 & 11.26\% & \code{undaktoven} & \code{INDACTOVEN} & 2 \\
Single Letter Phonetic Confusion & 5,983 & 9.03\% & \code{viratinantiseptic} & \code{VIRADIN ANTISEPTIC} & 1 \\
Multi Char Phonetic Confusion & 4,991 & 7.53\% & \code{kaminomotohairgrowthacceleratoryiupgrade} & \code{KAMINOMOTO HAIR GROWTH ACCELERATOR II UPGRADE} & 1 \\
Single Char Deletion Position Weighted & 4,966 & 7.50\% & \code{mutelaskinfreshener} & \code{MUSTELA SKIN FRESHENER} & 1 \\
Wrong Vowels Relabelled By Actual Changes & 4,870 & 7.35\% & \code{cupralex} & \code{CIPRALEX} & 1 \\
Combined Errors Effective 3 & 4,346 & 6.56\% & \code{qiabtivia} & \code{GLAPTIVIA} & 3 \\
OCR Letter Digit Confusion & 3,998 & 6.03\% & \code{imipenemcilastatinkab1} & \code{IMIPENEM CILASTATIN KABI} & 1 \\
Ligature Confusion & 3,992 & 6.03\% & \code{touliflyfacal} & \code{TOULIFLY FACIAL} & 1 \\
Transposition Position Weighted & 3,000 & 4.53\% & \code{chaplha} & \code{CH ALPHA} & 2 \\
OCR + Error Effective 2 & 2,997 & 4.52\% & \code{vitam1nkmeniq7} & \code{VITAMIN K MENAQ7} & 2 \\
Single Char Insertion & 2,988 & 4.51\% & \code{calciumviit} & \code{CALCIUM VIT} & 1 \\
Double Letter Reduction Expansion & 2,987 & 4.51\% & \code{ssimvacor} & \code{SIMVACOR} & 1 \\
Speed Typing Errors & 2,981 & 4.50\% & \code{rngeracetateelnasr} & \code{RINGER ACETATE EL NASR} & 1 \\
Keyboard Adjacent Sampled & 1,924 & 2.90\% & \code{billakordhair} & \code{VILLAKORD HAIR} & 1 \\
Combined Errors Effective 4 + & 563 & 0.85\% & \code{myclhaeme} & \code{MEDAHEME} & 4 \\
Combined Errors Effective 1 & 496 & 0.75\% & \code{whivewello} & \code{WHITE WELLO} & 1 \\
Dangerous Name Pair Gold Closer & 158 & 0.24\% & \code{casty} & \code{CASTO} & 1 \\
Consonant Frame Wrong Vowels Heavy & 68 & 0.10\% & \code{nopyl} & \code{NOPEL} & 1 \\
Extreme Combined Effective 2 & 4 & 0.01\% & \code{vfnt} & \code{VFEND} & 2 \\
Extreme Combined Effective 3 & 3 & 0.00\% & \code{mokce} & \code{MEGACE} & 3 \\
OCR + Error Effective 1 & 2 & 0.00\% & \code{ope} & \code{OPIE} & 1 \\
}

\newcommand{\ErrorTypeRows}{%
Wrong 1 Vowel Changes & 2,299 & 3.47\% & \code{cupralex} & \code{CIPRALEX} \\
Wrong 2 Vowel Changes & 1,824 & 2.75\% & \code{tobrialyx} & \code{TOBRA ALEX} \\
Wrong 3 Vowel Changes & 747 & 1.13\% & \code{metuniemevarydayeasy} & \code{METANIUM EVERYDAY EASY} \\
Double First Letter & 645 & 0.97\% & \code{vvalproexcr} & \code{VALPROEX CR} \\
Stutter Prefix 2 & 617 & 0.93\% & \code{chchampix} & \code{CHAMPIX} \\
Repeat Last Letter & 576 & 0.87\% & \code{unasynn} & \code{UNASYN} \\
Stutter Prefix 3 & 576 & 0.87\% & \code{bosbostonximajointh} & \code{BOSTON XIMAJOINT H} \\
Skip Second Letter & 567 & 0.86\% & \code{rngeracetateelnasr} & \code{RINGER ACETATE EL NASR} \\
Deletion First Char Pos 0 & 496 & 0.75\% & \code{ustelagentlecleansing} & \code{MUSTELA GENTLE CLEANSING} \\
Vowel + Phonetic Variant 3  Effective Damerau 2 & 435 & 0.66\% & \code{undaktoven} & \code{INDACTOVEN} \\
Deletion + Phonetic Variant 1  Effective Damerau 2 & 424 & 0.64\% & \code{darbrim} & \code{DARAPRIM} \\
Deletion + Phonetic Variant 3  Effective Damerau 2 & 410 & 0.62\% & \code{isiskingerlmonfilterbags} & \code{ISIS GINGER LEMON FILTER BAGS} \\
Vowel + Phonetic Variant 0  Effective Damerau 2 & 404 & 0.61\% & \code{polunehair} & \code{BOLINE HAIR} \\
Vowel + Phonetic Variant 2  Effective Damerau 2 & 404 & 0.61\% & \code{muxidrix} & \code{MOXITRIX} \\
Visual + Phonetic Variant 1  Effective Damerau 2 & 387 & 0.58\% & \code{phoztec} & \code{PROSTEC} \\
Deletion + Phonetic Variant 2  Effective Damerau 2 & 384 & 0.58\% & \code{rotakrbamide} & \code{ROTA CARBAMIDE} \\
Vowel + Phonetic Variant 1  Effective Damerau 2 & 380 & 0.57\% & \code{antocholyst} & \code{ANDOCHOLEST} \\
Transposition + Deletion Variant 0  Effective Damerau 2 & 378 & 0.57\% & \code{igvarrock} & \code{GIVAROROCK} \\
Visual + Phonetic Variant 3  Effective Damerau 2 & 374 & 0.56\% & \code{carhitopalevodopaeva} & \code{CARBIDOPA LEVODOPA EVA} \\
Deletion + Phonetic Variant 0  Effective Damerau 2 & 370 & 0.56\% & \code{trbsiflo} & \code{TRYPSIFLO} \\
Transposition + Deletion Variant 1  Effective Damerau 2 & 357 & 0.54\% & \code{craviprss} & \code{CARVIPRESS} \\
Visual Visual Phonetic Variant 1  Effective Damerau 3 & 344 & 0.52\% & \code{dhlxy} & \code{TRIXY} \\
Visual + Phonetic Variant 2  Effective Damerau 2 & 343 & 0.52\% & \code{zynorent} & \code{ZYMOREND} \\
Visual Visual Phonetic Variant 2  Effective Damerau 3 & 329 & 0.50\% & \code{icurtra} & \code{ICANDRA} \\
Vowel Phonetic Deletion Variant 2  Effective Damerau 3 & 329 & 0.50\% & \code{arckhuma} & \code{ARCECHIMA} \\
}

\newcommand{\OverallRows}{%
Exact or prefix match & 3.0019\% & 3.0125\% & 3.0125\% & 3.0125\% & 0.030072 & 1.6 & 0.7079 & 0.8049 & 0.06 & \textbf{0.0000\%} \\
Exhaustive Levenshtein & 94.0489\% & 98.3956\% & 99.0250\% & 99.4416\% & 0.959861 & 0.0 & 0.7357 & 0.7990 & 17476.00 & \textbf{0.0000\%} \\
Jaro-Winkler & 89.5694\% & 96.2630\% & 97.5459\% & 98.4273\% & 0.925647 & \textbf{0.0} & 0.7249 & 0.7907 & 17476.00 & \textbf{0.0000\%} \\
Character 3-gram TF-IDF & 75.2071\% & 86.7154\% & 89.4230\% & 91.6613\% & 0.804208 & 64.0 & \textbf{0.5126} & \textbf{0.6210} & 921.27 & \textbf{0.0000\%} \\
RapidFuzz token ratio & 65.6716\% & 78.2498\% & 82.0774\% & 85.4717\% & 0.713468 & 0.4 & 8.1414 & 12.2144 & 17476.00 & \textbf{0.0000\%} \\
Phonetic baseline & 69.6560\% & 84.9223\% & 89.0396\% & 92.0703\% & 0.764241 & 0.2 & 0.6924 & 0.7842 & 17475.74 & \textbf{0.0000\%} \\
Algorithm 1, current app & 83.8583\% & 89.7188\% & 91.2507\% & 92.2333\% & 0.863662 & 4722.1 & 1.2253 & 4.5748 & 35.91 & \textbf{0.0000\%} \\
Algorithm 2, external fast & 92.4053\% & 96.3325\% & 96.7324\% & 96.8788\% & 0.941355 & 4344.6 & 5.7286 & 21.9070 & 93.77 & 0.1675\% \\
Algorithm 3, rank fusion & 91.8363\% & 96.1830\% & 96.7309\% & 97.1203\% & 0.938051 & 12792.4 & 8.3464 & 24.4607 & 34.78 & \textbf{0.0000\%} \\
Algorithm 4, family rescue & 97.1852\% & 99.2423\% & 99.4250\% & 99.5155\% & 0.980944 & 8737.0 & 14.6203 & 31.5124 & 24.36 & \textbf{0.0000\%} \\
Algorithm 5, evidence-guided rescue & \textbf{98.2055\%} & \textbf{99.6619\%} & \textbf{99.8989\%} & \textbf{100.0000\%} & \textbf{0.988569} & 16750.6 & 52.1872 & 117.7635 & 24.54 & \textbf{0.0000\%} \\
}

\newcommand{\BehaviorRows}{%
Exact or prefix match & 3.22\% & 96.78\% & 0.0000\% & 0.06 \\
Exhaustive Levenshtein & 100.00\% & 0.00\% & 0.0000\% & 17476.00 \\
Jaro-Winkler & 100.00\% & 0.00\% & 0.0000\% & 17476.00 \\
Character 3-gram TF-IDF & 99.71\% & 0.29\% & 0.0000\% & 921.27 \\
RapidFuzz token ratio & 100.00\% & 0.00\% & 0.0000\% & 17476.00 \\
Phonetic baseline & 100.00\% & 0.00\% & 0.0000\% & 17475.74 \\
Algorithm 1, current app & 98.06\% & 1.94\% & 0.0000\% & 35.91 \\
Algorithm 2, external fast & 55.47\% & 0.15\% & 0.1675\% & 93.77 \\
Algorithm 3, rank fusion & 99.96\% & 0.04\% & 0.0000\% & 34.78 \\
Algorithm 4, family rescue & 99.95\% & 0.05\% & 0.0000\% & 24.36 \\
Algorithm 5, evidence-guided rescue & 100.00\% & 0.00\% & 0.0000\% & 24.54 \\
}

\newcommand{\PairedRows}{%
Exact or prefix match & +95.2035\pp & 63,079 & 0 & $<10^{-300}$ & +96.9875\pp & $<10^{-300}$ \\
Exhaustive Levenshtein & +4.1565\pp & 2,923 & 169 & $<10^{-300}$ & +0.5584\pp & $8.32\times10^{-112}$ \\
Jaro-Winkler & +8.6361\pp & 5,870 & 148 & $<10^{-300}$ & +1.5727\pp & $4.24\times10^{-314}$ \\
Character 3-gram TF-IDF & +22.9983\pp & 15,304 & 66 & $<10^{-300}$ & +8.3387\pp & $<10^{-300}$ \\
RapidFuzz token ratio & +32.5339\pp & 21,691 & 135 & $<10^{-300}$ & +14.5283\pp & $<10^{-300}$ \\
Phonetic baseline & +28.5494\pp & 19,065 & 149 & $<10^{-300}$ & +7.9297\pp & $<10^{-300}$ \\
Algorithm 1, current app & +14.3472\pp & 9,644 & 138 & $<10^{-300}$ & +7.7667\pp & $<10^{-300}$ \\
Algorithm 2, external fast & +5.8001\pp & 4,046 & 203 & $<10^{-300}$ & +3.1212\pp & $<10^{-300}$ \\
Algorithm 3, rank fusion & +6.3691\pp & 4,417 & 197 & $<10^{-300}$ & +2.8797\pp & $<10^{-300}$ \\
Algorithm 4, family rescue & +1.0203\pp & 688 & 12 & $1.02\times10^{-185}$ & +0.4845\pp & $4.68\times10^{-97}$ \\
}

\newcommand{\SplitRows}{%
Exact or prefix match & 2.95\% & 3.21\% & 2.96\% & 3.22\% \\
Exhaustive Levenshtein & 94.17\% & 93.57\% & 99.45\% & 99.42\% \\
Jaro-Winkler & 89.62\% & 89.38\% & 98.43\% & 98.43\% \\
Character 3-gram TF-IDF & 75.22\% & 75.15\% & 91.78\% & 91.18\% \\
RapidFuzz token ratio & 65.95\% & 64.56\% & 85.43\% & 85.63\% \\
Phonetic baseline & 69.52\% & 70.18\% & 92.08\% & 92.04\% \\
Algorithm 1, current app & 83.97\% & 83.42\% & 92.30\% & 91.97\% \\
Algorithm 2, external fast & 92.51\% & 91.99\% & 96.96\% & 96.56\% \\
Algorithm 3, rank fusion & 91.94\% & 91.44\% & 97.22\% & 96.73\% \\
Algorithm 4, family rescue & 97.24\% & 96.96\% & 99.53\% & 99.47\% \\
Algorithm 5, evidence-guided rescue & 98.22\% & 98.16\% & 100.00\% & 100.00\% \\
}

\newcommand{\DistanceHOneRows}{%
Exact or prefix match & 3.00\% & 5.14\% & 0.00\% & 0.00\% & 0.00\% & 0.00\% \\
Exhaustive Levenshtein & 94.05\% & 99.04\% & 91.89\% & 76.10\% & 43.47\% & 64.54\% \\
Jaro-Winkler & 89.57\% & 97.35\% & 84.24\% & 64.48\% & 47.03\% & 60.70\% \\
Character 3-gram TF-IDF & 75.21\% & 87.96\% & 62.44\% & 45.12\% & 21.14\% & 34.50\% \\
RapidFuzz token ratio & 65.67\% & 71.22\% & 62.33\% & 47.44\% & 29.22\% & 30.35\% \\
Phonetic baseline & 69.66\% & 72.60\% & 68.00\% & 60.64\% & 33.97\% & 51.76\% \\
Algorithm 1, current app & 83.86\% & 92.41\% & 80.86\% & 51.41\% & 13.06\% & 2.56\% \\
Algorithm 2, external fast & 92.41\% & 96.94\% & 90.66\% & 76.95\% & 37.77\% & 45.69\% \\
Algorithm 3, rank fusion & 91.84\% & 96.40\% & 90.36\% & 75.37\% & 36.82\% & 45.05\% \\
Algorithm 4, family rescue & 97.19\% & 99.55\% & 96.59\% & 88.29\% & 59.86\% & 81.79\% \\
Algorithm 5, evidence-guided rescue & \textbf{98.21\%} & \textbf{99.69\%} & \textbf{97.87\%} & \textbf{91.73\%} & \textbf{80.76\%} & \textbf{96.17\%} \\
}

\newcommand{\DistanceHTwentyRows}{%
Exact or prefix match & 3.01\% & 5.16\% & 0.00\% & 0.00\% & 0.00\% & 0.00\% \\
Exhaustive Levenshtein & 99.44\% & \textbf{100.00\%} & 99.98\% & 96.84\% & 73.87\% & 85.30\% \\
Jaro-Winkler & 98.43\% & 99.97\% & 98.62\% & 90.23\% & 83.37\% & 88.82\% \\
Character 3-gram TF-IDF & 91.66\% & 98.49\% & 86.82\% & 71.26\% & 41.09\% & 60.06\% \\
RapidFuzz token ratio & 85.47\% & 87.66\% & 85.00\% & 76.37\% & 66.51\% & 63.58\% \\
Phonetic baseline & 92.07\% & 94.08\% & 91.60\% & 84.53\% & 60.81\% & 75.40\% \\
Algorithm 1, current app & 92.23\% & 98.05\% & 94.27\% & 61.00\% & 15.68\% & 4.15\% \\
Algorithm 2, external fast & 96.88\% & 99.87\% & 96.91\% & 84.78\% & 46.79\% & 48.56\% \\
Algorithm 3, rank fusion & 97.12\% & 99.75\% & 97.41\% & 86.34\% & 46.79\% & 48.56\% \\
Algorithm 4, family rescue & 99.52\% & 99.90\% & 99.74\% & 98.15\% & 81.95\% & 89.78\% \\
Algorithm 5, evidence-guided rescue & \textbf{100.00\%} & \textbf{100.00\%} & \textbf{100.00\%} & \textbf{100.00\%} & \textbf{100.00\%} & \textbf{100.00\%} \\
}

\newcommand{\OperationRows}{%
Exact or prefix match & 0.0\% & 0.0\% & 0.1\% & 0.0\% & 48.9\% & 0.0\% & 0.0\% & 0.0\% & 0.2\% \\
Exhaustive Levenshtein & 84.2\% & 98.4\% & 98.9\% & 95.1\% & 96.4\% & 94.0\% & 99.8\% & 98.2\% & 99.0\% \\
Jaro-Winkler & 75.7\% & 95.3\% & 96.3\% & 79.2\% & 98.7\% & 99.4\% & 99.5\% & 92.4\% & 98.9\% \\
Character 3-gram TF-IDF & 51.2\% & 86.5\% & 84.4\% & 57.2\% & 93.5\% & 75.2\% & 97.2\% & 78.2\% & 99.4\% \\
RapidFuzz token ratio & 56.9\% & 67.5\% & 67.2\% & 64.5\% & 74.3\% & 66.0\% & 80.1\% & 74.6\% & 72.2\% \\
Phonetic baseline & 66.1\% & 64.3\% & 83.1\% & 78.3\% & 47.4\% & 81.2\% & 74.7\% & 47.5\% & 79.0\% \\
Algorithm 1, current app & 68.9\% & 84.9\% & 95.4\% & 96.3\% & 81.3\% & 94.9\% & 91.6\% & 75.7\% & 93.6\% \\
Algorithm 2, external fast & 82.9\% & 93.1\% & 97.9\% & 97.5\% & 97.8\% & 98.8\% & 98.4\% & 86.8\% & 99.5\% \\
Algorithm 3, rank fusion & 82.2\% & 93.0\% & 98.0\% & 97.4\% & 92.6\% & 98.1\% & 97.8\% & 87.3\% & 99.4\% \\
Algorithm 4, family rescue & 92.5\% & 98.5\% & 99.4\% & 98.5\% & 99.5\% & 99.5\% & 99.6\% & 98.1\% & 99.6\% \\
Algorithm 5, evidence-guided rescue & 95.0\% & 99.2\% & 99.8\% & 98.8\% & 99.6\% & 100.0\% & 99.6\% & 98.5\% & 100.0\% \\
}

\newcommand{\BigramRows}{%
Exact or prefix match & 3.00\% & 0.00\% & 0.00\% & 1.28\% & 3.28\% \\
Exhaustive Levenshtein & 94.05\% & 9.74\% & 33.22\% & 71.24\% & 98.41\% \\
Jaro-Winkler & 89.57\% & 18.05\% & 29.90\% & 60.50\% & 94.55\% \\
Character 3-gram TF-IDF & 75.21\% & 0.00\% & 0.00\% & 17.04\% & 83.90\% \\
RapidFuzz token ratio & 65.67\% & 16.05\% & 25.35\% & 59.50\% & 67.48\% \\
Phonetic baseline & 69.66\% & 21.20\% & 22.38\% & 34.92\% & 74.92\% \\
Algorithm 1, current app & 83.86\% & 26.93\% & 30.86\% & 57.77\% & 88.28\% \\
Algorithm 2, external fast & 92.41\% & 37.82\% & 39.60\% & 70.52\% & 96.32\% \\
Algorithm 3, rank fusion & 91.84\% & 32.38\% & 36.71\% & 69.77\% & 95.84\% \\
Algorithm 4, family rescue & 97.19\% & 40.11\% & 61.80\% & 87.61\% & 99.34\% \\
Algorithm 5, evidence-guided rescue & 98.21\% & 57.02\% & 74.30\% & 91.89\% & 99.66\% \\
}

\newcommand{\CohortRows}{%
Exact or prefix match & 64,664 & 3.08\% & 3.09\% & 1,470 & 0.00\% & 0.00\% & 123 & 0.00\% & 0.00\% \\
Exhaustive Levenshtein & 64,664 & 95.83\% & 99.99\% & 1,470 & 23.47\% & 81.50\% & 123 & 0.81\% & 26.83\% \\
Jaro-Winkler & 64,664 & 91.10\% & 99.25\% & 1,470 & 29.39\% & 66.94\% & 123 & 3.25\% & 40.65\% \\
Character 3-gram TF-IDF & 64,664 & 76.76\% & 93.30\% & 1,470 & 12.38\% & 26.60\% & 123 & 7.32\% & 8.13\% \\
RapidFuzz token ratio & 64,664 & 66.72\% & 86.10\% & 1,470 & 24.63\% & 62.38\% & 123 & 6.50\% & 32.52\% \\
Phonetic baseline & 64,664 & 70.94\% & 93.02\% & 1,470 & 19.12\% & 56.12\% & 123 & 0.81\% & 20.33\% \\
Algorithm 1, current app & 64,664 & 85.41\% & 93.66\% & 1,470 & 22.65\% & 36.80\% & 123 & 1.63\% & 5.69\% \\
Algorithm 2, external fast & 64,664 & 93.81\% & 98.02\% & 1,470 & 38.16\% & 53.95\% & 123 & 4.07\% & 8.94\% \\
Algorithm 3, rank fusion & 64,664 & 93.28\% & 98.24\% & 1,470 & 35.71\% & 55.03\% & 123 & 4.88\% & 12.20\% \\
Algorithm 4, family rescue & 64,664 & 98.35\% & 99.87\% & 1,470 & 52.99\% & 89.46\% & 123 & 13.82\% & 31.71\% \\
Algorithm 5, evidence-guided rescue & 64,664 & 98.93\% & 100.00\% & 1,470 & 69.05\% & 100.00\% & 123 & 68.29\% & 100.00\% \\
}

\newcommand{\ExtremeDistributionRows}{%
Effective edits & 2 & 12 & 9.76\% & 100.00\% & 100.00\% \\
Effective edits & 3 & 9 & 7.32\% & 100.00\% & 100.00\% \\
Effective edits & 4 & 58 & 47.15\% & 44.83\% & 100.00\% \\
Effective edits & 5 & 44 & 35.77\% & 84.09\% & 100.00\% \\
Operation family & Mixed Operations & 104 & 84.55\% & 62.50\% & 100.00\% \\
Operation family & Other & 9 & 7.32\% & 100.00\% & 100.00\% \\
Operation family & Phonetic & 1 & 0.81\% & 100.00\% & 100.00\% \\
Operation family & Transposition & 5 & 4.07\% & 100.00\% & 100.00\% \\
Operation family & Visual Or OCR & 4 & 3.25\% & 100.00\% & 100.00\% \\
Shared bigrams & 0 Shared Bigrams & 41 & 33.33\% & 51.22\% & 100.00\% \\
Shared bigrams & 1 Shared Bigram & 51 & 41.46\% & 72.55\% & 100.00\% \\
Shared bigrams & 2 3 Shared Bigrams & 31 & 25.20\% & 83.87\% & 100.00\% \\
Difficulty & Easy & 2 & 1.63\% & 100.00\% & 100.00\% \\
Difficulty & Extreme & 61 & 49.59\% & 65.57\% & 100.00\% \\
Difficulty & Hard & 54 & 43.90\% & 66.67\% & 100.00\% \\
Difficulty & Medium & 6 & 4.88\% & 100.00\% & 100.00\% \\
}

\newcommand{\ExtremeExampleRows}{%
\code{pryozlarn} & \code{PRAZOLAM} & 0.625 & 2 & \code{PRAZOLAM} & 1 & Combined Errors Effective 4 + \\
\code{lvce} & \code{VOLCA} & 0.800 & 0 & \code{LEVCET} & 4 & Combined Errors Effective 3 \\
\code{milmila} & \code{MILA} & 0.750 & 3 & \code{MILA} & 1 & Speed Typing Errors \\
\code{rnol} & \code{MOL} & 0.667 & 1 & \code{MOL} & 1 & Ligature Confusion \\
\code{ivee} & \code{IVY} & 0.667 & 1 & \code{IVY} & 1 & Multi Char Phonetic Confusion \\
\code{pb5} & \code{BBS} & 0.667 & 0 & \code{BBS} & 1 & OCR + Error Effective 2 \\
\code{2ix} & \code{ZEX} & 0.667 & 0 & \code{ZEX} & 1 & OCR + Error Effective 2 \\
\code{nta} & \code{NAT} & 0.667 & 0 & \code{NAT} & 1 & Transposition Position Weighted \\
\code{alx} & \code{LAX} & 0.667 & 0 & \code{LAX} & 1 & Transposition Position Weighted \\
\code{ful} & \code{FLU} & 0.667 & 0 & \code{FLU} & 1 & Transposition Position Weighted \\
}

\newcommand{\ClassicalCategoryRows}{%
Single Letter Visual Confusion & 7,841 & 0.0\% & 100.0\% & 100.0\% & 98.7\% & 86.4\% & 91.7\% \\
Combined Errors Effective 2 & 7,463 & 0.0\% & 99.5\% & 98.1\% & 75.3\% & 85.4\% & 93.6\% \\
Single Letter Phonetic Confusion & 5,983 & 0.0\% & 100.0\% & 99.9\% & 98.1\% & 85.5\% & 100.0\% \\
Multi Char Phonetic Confusion & 4,991 & 0.3\% & 100.0\% & 99.9\% & 98.6\% & 86.8\% & 95.4\% \\
Single Char Deletion Position Weighted & 4,966 & 39.8\% & 100.0\% & 100.0\% & 99.2\% & 87.9\% & 82.6\% \\
Wrong Vowels Relabelled By Actual Changes & 4,870 & 0.0\% & 99.9\% & 97.0\% & 84.2\% & 85.2\% & 100.0\% \\
Combined Errors Effective 3 & 4,346 & 0.0\% & 94.3\% & 86.1\% & 65.3\% & 73.9\% & 77.8\% \\
OCR Letter Digit Confusion & 3,998 & 0.0\% & 100.0\% & 100.0\% & 98.9\% & 87.3\% & 93.1\% \\
Ligature Confusion & 3,992 & 0.0\% & 100.0\% & 99.4\% & 97.4\% & 84.1\% & 78.7\% \\
Transposition Position Weighted & 3,000 & 0.0\% & 100.0\% & 99.8\% & 94.9\% & 85.1\% & 95.7\% \\
OCR + Error Effective 2 & 2,997 & 0.0\% & 100.0\% & 99.5\% & 88.0\% & 82.3\% & 97.9\% \\
Single Char Insertion & 2,988 & 0.0\% & 100.0\% & 100.0\% & 99.7\% & 90.7\% & 98.8\% \\
Double Letter Reduction Expansion & 2,987 & 0.3\% & 100.0\% & 100.0\% & 99.9\% & 88.5\% & 100.0\% \\
Speed Typing Errors & 2,981 & 0.0\% & 100.0\% & 100.0\% & 99.8\% & 90.7\% & 94.6\% \\
Keyboard Adjacent Sampled & 1,924 & 0.0\% & 100.0\% & 99.9\% & 94.4\% & 88.2\% & 80.9\% \\
Combined Errors Effective 4 + & 563 & 0.0\% & 86.0\% & 83.7\% & 54.9\% & 63.4\% & 68.7\% \\
Combined Errors Effective 1 & 496 & 0.4\% & 100.0\% & 100.0\% & 94.0\% & 91.3\% & 95.4\% \\
Dangerous Name Pair Gold Closer & 158 & 0.0\% & 100.0\% & 100.0\% & 80.4\% & 92.4\% & 97.5\% \\
Consonant Frame Wrong Vowels Heavy & 68 & 0.0\% & 100.0\% & 91.2\% & 66.2\% & 80.9\% & 100.0\% \\
Extreme Combined Effective 2 & 4 & 0.0\% & 100.0\% & 100.0\% & 50.0\% & 100.0\% & 100.0\% \\
Extreme Combined Effective 3 & 3 & 0.0\% & 100.0\% & 66.7\% & 66.7\% & 33.3\% & 100.0\% \\
OCR + Error Effective 1 & 2 & 0.0\% & 100.0\% & 100.0\% & 50.0\% & 100.0\% & 100.0\% \\
}

\newcommand{\CategoryAlgorithmRows}{%
Single Letter Visual Confusion & 7,841 & 88.2\% & 94.5\% & 94.4\% & 99.6\% & 99.7\% & 97.3\% & 99.9\% & 99.8\% & 100.0\% & 100.0\% \\
Combined Errors Effective 2 & 7,463 & 77.8\% & 87.1\% & 86.1\% & 95.3\% & 96.1\% & 92.0\% & 95.7\% & 96.0\% & 99.7\% & 100.0\% \\
Single Letter Phonetic Confusion & 5,983 & 98.8\% & 97.4\% & 98.2\% & 99.6\% & 99.8\% & 99.6\% & 99.9\% & 99.9\% & 99.9\% & 100.0\% \\
Multi Char Phonetic Confusion & 4,991 & 91.3\% & 98.5\% & 97.7\% & 99.2\% & 99.8\% & 97.7\% & 99.8\% & 99.7\% & 99.9\% & 100.0\% \\
Single Char Deletion Position Weighted & 4,966 & 86.7\% & 98.4\% & 94.8\% & 99.6\% & 99.7\% & 95.7\% & 99.9\% & 99.8\% & 100.0\% & 100.0\% \\
Wrong Vowels Relabelled By Actual Changes & 4,870 & 96.3\% & 97.5\% & 97.5\% & 98.5\% & 98.8\% & 98.7\% & 99.0\% & 99.1\% & 99.5\% & 100.0\% \\
Combined Errors Effective 3 & 4,346 & 34.1\% & 64.4\% & 62.6\% & 81.5\% & 86.7\% & 45.5\% & 75.6\% & 77.9\% & 96.4\% & 100.0\% \\
OCR Letter Digit Confusion & 3,998 & 94.2\% & 96.9\% & 96.7\% & 99.5\% & 99.8\% & 98.9\% & 99.8\% & 99.9\% & 99.8\% & 100.0\% \\
Ligature Confusion & 3,992 & 69.1\% & 86.5\% & 86.5\% & 95.2\% & 97.7\% & 94.5\% & 97.4\% & 98.1\% & 99.9\% & 100.0\% \\
Transposition Position Weighted & 3,000 & 94.9\% & 98.8\% & 98.1\% & 99.5\% & 100.0\% & 98.4\% & 99.8\% & 99.8\% & 99.8\% & 100.0\% \\
OCR + Error Effective 2 & 2,997 & 81.9\% & 91.4\% & 92.5\% & 97.2\% & 99.2\% & 92.1\% & 94.3\% & 96.0\% & 99.5\% & 100.0\% \\
Single Char Insertion & 2,988 & 91.8\% & 98.7\% & 98.0\% & 99.7\% & 99.7\% & 99.0\% & 100.0\% & 99.8\% & 100.0\% & 100.0\% \\
Double Letter Reduction Expansion & 2,987 & 98.8\% & 99.9\% & 99.8\% & 99.9\% & 100.0\% & 99.9\% & 99.9\% & 99.9\% & 99.9\% & 100.0\% \\
Speed Typing Errors & 2,981 & 86.6\% & 98.6\% & 97.4\% & 99.3\% & 99.9\% & 90.2\% & 99.8\% & 99.8\% & 100.0\% & 100.0\% \\
Keyboard Adjacent Sampled & 1,924 & 75.7\% & 86.8\% & 87.2\% & 98.1\% & 98.5\% & 92.7\% & 99.6\% & 98.7\% & 99.6\% & 100.0\% \\
Combined Errors Effective 4 + & 563 & 2.0\% & 39.3\% & 38.5\% & 73.5\% & 94.1\% & 3.0\% & 45.8\% & 45.6\% & 88.1\% & 100.0\% \\
Combined Errors Effective 1 & 496 & 91.5\% & 96.4\% & 94.8\% & 99.8\% & 99.8\% & 98.6\% & 100.0\% & 100.0\% & 100.0\% & 100.0\% \\
Dangerous Name Pair Gold Closer & 158 & 82.9\% & 87.3\% & 85.4\% & 97.5\% & 98.1\% & 91.8\% & 98.7\% & 98.1\% & 100.0\% & 100.0\% \\
Consonant Frame Wrong Vowels Heavy & 68 & 91.2\% & 97.1\% & 95.6\% & 97.1\% & 97.1\% & 97.1\% & 100.0\% & 100.0\% & 100.0\% & 100.0\% \\
Extreme Combined Effective 2 & 4 & 50.0\% & 75.0\% & 75.0\% & 75.0\% & 75.0\% & 75.0\% & 100.0\% & 75.0\% & 100.0\% & 100.0\% \\
Extreme Combined Effective 3 & 3 & 66.7\% & 66.7\% & 66.7\% & 66.7\% & 66.7\% & 66.7\% & 66.7\% & 66.7\% & 100.0\% & 100.0\% \\
OCR + Error Effective 1 & 2 & 50.0\% & 100.0\% & 50.0\% & 100.0\% & 100.0\% & 50.0\% & 100.0\% & 100.0\% & 100.0\% & 100.0\% \\
}

\newcommand{\AFiveImpactRows}{%
Combined Errors Effective 3 & 4,346 & 86.75\% & 100.00\% & 576 & 13.25\% & 48.44\% \\
Combined Errors Effective 2 & 7,463 & 96.07\% & 100.00\% & 293 & 3.93\% & 24.64\% \\
Ligature Confusion & 3,992 & 97.72\% & 100.00\% & 91 & 2.28\% & 7.65\% \\
Wrong Vowels Relabelled By Actual Changes & 4,870 & 98.79\% & 100.00\% & 59 & 1.21\% & 4.96\% \\
Combined Errors Effective 4 + & 563 & 94.14\% & 100.00\% & 33 & 5.86\% & 2.78\% \\
Keyboard Adjacent Sampled & 1,924 & 98.54\% & 100.00\% & 28 & 1.46\% & 2.35\% \\
Single Letter Visual Confusion & 7,841 & 99.69\% & 100.00\% & 24 & 0.31\% & 2.02\% \\
OCR + Error Effective 2 & 2,997 & 99.23\% & 100.00\% & 23 & 0.77\% & 1.93\% \\
Single Char Deletion Position Weighted & 4,966 & 99.66\% & 100.00\% & 17 & 0.34\% & 1.43\% \\
Single Letter Phonetic Confusion & 5,983 & 99.80\% & 100.00\% & 12 & 0.20\% & 1.01\% \\
Multi Char Phonetic Confusion & 4,991 & 99.80\% & 100.00\% & 10 & 0.20\% & 0.84\% \\
OCR Letter Digit Confusion & 3,998 & 99.77\% & 100.00\% & 9 & 0.23\% & 0.76\% \\
Single Char Insertion & 2,988 & 99.70\% & 100.00\% & 9 & 0.30\% & 0.76\% \\
Speed Typing Errors & 2,981 & 99.87\% & 100.00\% & 4 & 0.13\% & 0.34\% \\
Dangerous Name Pair Gold Closer & 158 & 98.10\% & 100.00\% & 3 & 1.90\% & 0.25\% \\
Consonant Frame Wrong Vowels Heavy & 68 & 97.06\% & 100.00\% & 2 & 2.94\% & 0.17\% \\
Combined Errors Effective 1 & 496 & 99.80\% & 100.00\% & 1 & 0.20\% & 0.08\% \\
Extreme Combined Effective 2 & 4 & 75.00\% & 100.00\% & 1 & 25.00\% & 0.08\% \\
Extreme Combined Effective 3 & 3 & 66.67\% & 100.00\% & 1 & 33.33\% & 0.08\% \\
Transposition Position Weighted & 3,000 & 100.00\% & 100.00\% & 0 & 0.00\% & 0.00\% \\
Double Letter Reduction Expansion & 2,987 & 100.00\% & 100.00\% & 0 & 0.00\% & 0.00\% \\
OCR + Error Effective 1 & 2 & 100.00\% & 100.00\% & 0 & 0.00\% & 0.00\% \\
}

\newcommand{\DangerRows}{%
Exact or prefix match & 158 & 0.00\% & 0.00\% & 0.00\% \\
Exhaustive Levenshtein & 158 & 86.71\% & 100.00\% & 0.00\% \\
Jaro-Winkler & 158 & 81.01\% & 100.00\% & 0.00\% \\
Character 3-gram TF-IDF & 158 & 35.44\% & 80.38\% & 0.00\% \\
RapidFuzz token ratio & 158 & 75.95\% & 92.41\% & 0.00\% \\
Phonetic baseline & 158 & 43.04\% & 97.47\% & 0.00\% \\
Algorithm 1, current app & 158 & 82.91\% & 91.77\% & 0.00\% \\
Algorithm 2, external fast & 158 & 87.34\% & 98.73\% & 0.63\% \\
Algorithm 3, rank fusion & 158 & 85.44\% & 98.10\% & 0.00\% \\
Algorithm 4, family rescue & 158 & 97.47\% & 100.00\% & 0.00\% \\
Algorithm 5, evidence-guided rescue & 158 & 98.10\% & 100.00\% & 0.00\% \\
}

\newcommand{\DangerExampleRows}{%
\code{casty} & \code{CASTO} & \code{CASHO} & 1 & 2 & Gold family is strictly closer, so one expected answer remains defensible. \\
\code{solvyn} & \code{SOLVIN} & \code{SOLVIA} & 1 & 2 & Gold family is strictly closer, so one expected answer remains defensible. \\
\code{levy} & \code{LEVA} & \code{LECA} & 1 & 2 & Gold family is strictly closer, so one expected answer remains defensible. \\
\code{vtop} & \code{VITOP} & \code{VITON} & 1 & 2 & Gold family is strictly closer, so one expected answer remains defensible. \\
\code{imix} & \code{E MIX} & \code{EMUX} & 1 & 2 & Gold family is strictly closer, so one expected answer remains defensible. \\
\code{mopram} & \code{MEPRAM} & \code{DEPRAM} & 1 & 2 & Gold family is strictly closer, so one expected answer remains defensible. \\
\code{lmex} & \code{LOMEX} & \code{TOMEX} & 1 & 2 & Gold family is strictly closer, so one expected answer remains defensible. \\
\code{rlly} & \code{ROLLY} & \code{JOLLY} & 1 & 2 & Gold family is strictly closer, so one expected answer remains defensible. \\
}

\newcommand{\EqualDistanceRows}{%
Exact or prefix match & 146 & 25 & 17.12\% & 7 & 18 \\
Exhaustive Levenshtein & 3,943 & 2,957 & 74.99\% & 2,933 & 24 \\
Jaro-Winkler & 6,911 & 1,754 & 25.38\% & 1,409 & 345 \\
Character 3-gram TF-IDF & 16,238 & 1,317 & 8.11\% & 668 & 649 \\
RapidFuzz token ratio & 22,745 & 2,131 & 9.37\% & 1,716 & 415 \\
Phonetic baseline & 20,104 & 324 & 1.61\% & 201 & 123 \\
Algorithm 1, current app & 9,412 & 566 & 6.01\% & 304 & 262 \\
Algorithm 2, external fast & 4,931 & 867 & 17.58\% & 490 & 377 \\
Algorithm 3, rank fusion & 5,383 & 841 & 15.62\% & 514 & 327 \\
Algorithm 4, family rescue & 1,834 & 934 & 50.93\% & 855 & 79 \\
Algorithm 5, evidence-guided rescue & 1,189 & 563 & 47.35\% & 563 & 0 \\
}

\newcommand{\EqualDistanceExampleRows}{%
\code{coex} & \code{COPEX} & \code{COGEX} & 1 & 2 & Combined Errors Effective 1 & ranking \\
\code{mepo} & \code{MEBO} & \code{MEPRO} & 1 & 2 & Dangerous Name Pair Gold Closer & ranking \\
\code{roaix} & \code{ROSIX} & \code{ROMIX} & 1 & 2 & Keyboard Adjacent Sampled & ranking \\
\code{ultramdx} & \code{ULTRAMOX} & \code{ULTRA MIX} & 1 & 2 & OCR Letter Digit Confusion & ranking \\
\code{oriana} & \code{SORIANA} & \code{ORIANNA} & 1 & 2 & Single Char Deletion Position Weighted & ranking \\
\code{bnohair} & \code{NO HAIR} & \code{BIO HAIR} & 1 & 3 & Single Char Insertion & ranking \\
\code{opental} & \code{PENTAL} & \code{ORENTAL} & 1 & 2 & Single Char Insertion & ranking \\
\code{iotal} & \code{OTAL} & \code{TOTAL} & 1 & 3 & Single Char Insertion & ranking \\
\code{rollyt} & \code{ROLYT} & \code{ROLLY} & 1 & 2 & Single Char Insertion & ranking \\
\code{varioc} & \code{VARIC} & \code{VARIO} & 1 & 2 & Single Char Insertion & ranking \\
}

\newcommand{\AFiveGainRows}{%
Exact or prefix match & \code{claremnumhairconditioner} & \code{DERMENUM HAIR CONDITIONER} & \code{} & $>20$ & A5 ranks the verified family first at 5 edits; Combined Errors Effective 4 +. \\
Exhaustive Levenshtein & \code{rivisrpire} & \code{RIVAROSPIRE} & \code{IVY SPIR} & 3 & A5 ranks the verified family first at 4 edits; Combined Errors Effective 3. \\
Jaro-Winkler & \code{erxa} & \code{EURAX} & \code{SEROXAT} & 12 & A5 ranks the verified family first at 3 edits; Combined Errors Effective 2. \\
Character 3-gram TF-IDF & \code{delusav} & \code{DOLOSEV} & \code{CALUSA} & $>20$ & A5 ranks the verified family first at 3 edits; Consonant Frame Wrong Vowels Heavy. \\
RapidFuzz token ratio & \code{forsurehaircondishuning} & \code{FORSURE HAIR CONDITIONING} & \code{MOSHING HAIR} & $>20$ & A5 ranks the verified family first at 3 edits; Multi Char Phonetic Confusion. \\
Phonetic baseline & \code{vitvitamine} & \code{VITAMIN E} & \code{BETAVIDONE} & $>20$ & A5 ranks the verified family first at 3 edits; Speed Typing Errors. \\
Algorithm 1, current app & \code{haurvevu} & \code{HAIR VIVE} & \code{HAIR BACK} & 4 & A5 ranks the verified family first at 3 edits; Wrong Vowels Relabelled By Actual Changes. \\
Algorithm 2, external fast & \code{vitrolieal} & \code{VITROHEAL} & \code{OVITRELLE} & 2 & A5 ranks the verified family first at 2 edits; Ligature Confusion. \\
Algorithm 3, rank fusion & \code{xoal} & \code{XOLA} & \code{COOL} & 5 & A5 ranks the verified family first at 2 edits; Dangerous Name Pair Gold Closer. \\
Algorithm 4, family rescue & \code{0pinov} & \code{OPONOV} & \code{EPINOR} & 2 & A5 ranks the verified family first at 2 edits; OCR + Error Effective 2. \\
}

\newcommand{\FailureExampleRows}{%
Exact or prefix match & ranking & \code{cefzi} & \code{CEFZIM} & \code{CEFZIL} & 2 & 1 & Single Char Deletion Position Weighted \\
Exact or prefix match & retrieval & \code{zrtil} & \code{ROZTEL} & \code{} & $>20$ & 4 & Combined Errors Effective 3 \\
Exhaustive Levenshtein & ranking & \code{tarivn} & \code{TARIVAN} & \code{TARION} & 2 & 1 & Combined Errors Effective 1 \\
Exhaustive Levenshtein & retrieval & \code{zrtil} & \code{ROZTEL} & \code{ACNIL} & $>20$ & 4 & Combined Errors Effective 3 \\
Jaro-Winkler & ranking & \code{easyrezt} & \code{EASYREST} & \code{EASYRECT} & 2 & 1 & Combined Errors Effective 1 \\
Jaro-Winkler & retrieval & \code{liasat} & \code{HISTA} & \code{LIPOVAST} & $>20$ & 4 & Combined Errors Effective 3 \\
Character 3-gram TF-IDF & ranking & \code{pschodal} & \code{PSYCHODAL} & \code{ARCHODAL} & 2 & 1 & Combined Errors Effective 1 \\
Character 3-gram TF-IDF & retrieval & \code{zrtil} & \code{ROZTEL} & \code{CORTILON} & $>20$ & 4 & Combined Errors Effective 3 \\
RapidFuzz token ratio & ranking & \code{easyrezt} & \code{EASYREST} & \code{EASYRECT} & 2 & 1 & Combined Errors Effective 1 \\
RapidFuzz token ratio & retrieval & \code{zrtil} & \code{ROZTEL} & \code{EZAPRIL} & $>20$ & 4 & Combined Errors Effective 3 \\
Phonetic baseline & ranking & \code{hisvonic} & \code{HISTONIC} & \code{HEPAFENCE} & 2 & 1 & Combined Errors Effective 1 \\
Phonetic baseline & retrieval & \code{zrtil} & \code{ROZTEL} & \code{APIFORTYL} & $>20$ & 4 & Combined Errors Effective 3 \\
Algorithm 1, current app & ranking & \code{pflox} & \code{PEFLOX} & \code{PELOX} & 2 & 1 & Combined Errors Effective 1 \\
Algorithm 1, current app & retrieval & \code{zrtil} & \code{ROZTEL} & \code{ZEROL} & $>20$ & 4 & Combined Errors Effective 3 \\
Algorithm 2, external fast & ranking & \code{exludon} & \code{EXLUTON} & \code{OXELADINE} & 2 & 1 & Combined Errors Effective 1 \\
Algorithm 2, external fast & retrieval & \code{zrtil} & \code{ROZTEL} & \code{SUREPTIL} & $>20$ & 4 & Combined Errors Effective 3 \\
Algorithm 3, rank fusion & ranking & \code{exludon} & \code{EXLUTON} & \code{OXELADINE} & 2 & 1 & Combined Errors Effective 1 \\
Algorithm 3, rank fusion & retrieval & \code{zrtil} & \code{ROZTEL} & \code{SARTALESS} & $>20$ & 4 & Combined Errors Effective 3 \\
Algorithm 4, family rescue & ranking & \code{vidop} & \code{VITOP} & \code{VIDROP} & 2 & 1 & Dangerous Name Pair Gold Closer \\
Algorithm 4, family rescue & retrieval & \code{zrtil} & \code{ROZTEL} & \code{SORTI} & $>20$ & 4 & Combined Errors Effective 3 \\
Algorithm 5, evidence-guided rescue & ranking & \code{rocare} & \code{RECARE} & \code{PROCARE} & 3 & 1 & Dangerous Name Pair Gold Closer \\
}

\newcommand{\ScoreReconciliationRows}{%
Correct at rank 1 & 65,068 & 98.2055\% \\
Expected at ranks 2--5 & 965 & 1.4564\% \\
Expected at ranks 6--10 & 157 & 0.2370\% \\
Expected at ranks 11--20 & 67 & 0.1011\% \\
Expected outside top 20 & 0 & 0.0000\% \\
}

\newcommand{\FailureValidityRows}{%
A5 rows & 66,257 & Must equal the clean-core denominator. \\
Unique A5 case IDs & 66,257 & One result row per test case. \\
Recomputed Hit@1 mismatches & 0 & Top name was compared again with the expected family key. \\
Recomputed Hit@20 mismatches & 0 & Stored rank and Hit@20 flag must agree. \\
Levenshtein mismatches & 0 & Every stored distance was recalculated from query and target. \\
Unrestricted Damerau mismatches & 0 & The stored effective-Damerau field was recalculated with its own distance definition. \\
Stored Damerau versus A5 OSA differences & 88 & These are expected definition differences, not Hit@k scoring errors. \\
Exact different-family queries & 0 & A positive value would be an unfair real-name collision. \\
Unique-nearest failures rerun & 14 & Current A5 was executed again with candidate scores exposed. \\
Unique-nearest rerun mismatches & 0 & Fresh top name and expected rank must match the benchmark. \\
A5/OSA distance test pairs & 10,000 & Random string pairs used to verify the distance implementation. \\
A5/OSA implementation mismatches & 0 & A positive value would invalidate OSA geometry labels. \\
}

\newcommand{\FreshRerunCaption}{%

}

\newcommand{\FreshRerunAuditText}{%

}

\newcommand{\FamilyDenominatorRows}{%
Pair micro-average & 66,257 & 98.2055\% & 100.0000\% & Every distinct query-target pair receives one vote. \\
Family macro-average & 16,845 & 97.5120\% & 100.0000\% & Every verified family receives one vote after averaging its cases. \\
}

\newcommand{\FailureConcentrationRows}{%
Families with at least one Hit@1 failure & 1,070 & 6.35\% \\
Families with at least one Hit@20 failure & 0 & 0.00\% \\
Top ten families' share of Hit@1 failures & 27 & 2.27\% \\
}

\newcommand{\CatalogGeometryRows}{%
Gold is the unique nearest family & 61,228 & 92.41\% & 99.9771\% & 100.0000\% & 14 \\
Gold ties with catalog neighbors & 4,210 & 6.35\% & 86.4133\% & 100.0000\% & 572 \\
Another catalog family is closer & 819 & 1.24\% & 26.3736\% & 100.0000\% & 603 \\
}

\newcommand{\FailureRootRows}{%
Gold is the unique nearest family & 14 & 1.18\% & 14 & 0 \\
Gold ties with catalog neighbors & 572 & 48.11\% & 572 & 0 \\
Another catalog family is closer & 603 & 50.71\% & 603 & 0 \\
}

\newcommand{\FailureDistanceDefinitionRows}{%
Gold is the unique nearest family & 15 & 1.26\% & 14 & 1.18\% \\
Gold ties with catalog neighbors & 566 & 47.60\% & 572 & 48.11\% \\
Another catalog family is closer & 608 & 51.14\% & 603 & 50.71\% \\
}

\newcommand{\SystemGeometryRows}{%
Exact or prefix match & 3.01\% & 3.01\% & 3.40\% & 3.56\% & 0.00\% & 0.00\% \\
Exhaustive Levenshtein & 99.32\% & 100.00\% & 35.56\% & 97.74\% & 0.37\% & 66.79\% \\
Jaro-Winkler & 94.42\% & 99.73\% & 34.89\% & 88.08\% & 8.06\% & 54.46\% \\
Character 3-gram TF-IDF & 80.25\% & 95.46\% & 15.89\% & 49.88\% & 2.93\% & 22.34\% \\
RapidFuzz token ratio & 68.74\% & 85.77\% & 33.11\% & 87.67\% & 3.66\% & 51.53\% \\
Phonetic baseline & 73.13\% & 93.85\% & 29.43\% & 73.47\% & 16.85\% & 54.46\% \\
Algorithm 1, current app & 87.70\% & 94.79\% & 41.33\% & 66.01\% & 15.14\% & 36.14\% \\
Algorithm 2, external fast & 95.84\% & 98.69\% & 56.10\% & 79.52\% & 21.98\% & 50.55\% \\
Algorithm 3, rank fusion & 95.39\% & 98.86\% & 53.97\% & 80.83\% & 20.76\% & 50.43\% \\
Algorithm 4, family rescue & 99.85\% & 99.93\% & 75.80\% & 97.24\% & 7.69\% & 80.34\% \\
Algorithm 5, evidence-guided rescue & 99.98\% & 100.00\% & 86.41\% & 100.00\% & 26.37\% & 100.00\% \\
}

\newcommand{\UniqueFailureMechanismRows}{%
Combined-score inversion & 14 & 1.18\% & Expected family is present, uniquely nearest, but has a lower combined score. \\
Bounded correction override & 0 & 0.00\% & A correction moved a lower-score family above the uniquely nearest target. \\
No candidates returned & 0 & 0.00\% & Candidate generation returned an empty list. \\
Target outside top 20 & 0 & 0.00\% & Other candidates were returned, but the uniquely nearest target was absent. \\
}

\newcommand{\UniqueFailurePatternRows}{%
Expected already in top 20 & 14 & 8 & 0 & 0 & Reranking problem: 14 combined-score inversions and 0 bounded-correction overrides. \\
Expected outside top 20 & 0 & 0 & 0 & 0 & No remaining candidate-generation failures; every verified family is in the top 20. \\
}

\newcommand{\FailureRescueRows}{%
Algorithm 2, external fast & 203 & 17.07\% & 621 & 52.23\% \\
Algorithm 3, rank fusion & 197 & 16.57\% & 626 & 52.65\% \\
Exhaustive Levenshtein & 169 & 14.21\% & 982 & 82.59\% \\
Phonetic baseline & 149 & 12.53\% & 642 & 53.99\% \\
Jaro-Winkler & 148 & 12.45\% & 764 & 64.26\% \\
Algorithm 1, current app & 138 & 11.61\% & 429 & 36.08\% \\
RapidFuzz token ratio & 135 & 11.35\% & 776 & 65.26\% \\
Character 3-gram TF-IDF & 66 & 5.55\% & 368 & 30.95\% \\
Algorithm 4, family rescue & 12 & 1.01\% & 1,020 & 85.79\% \\
Exact or prefix match & 0 & 0.00\% & 5 & 0.42\% \\
}

\newcommand{\AllSystemMissRows}{%
Total & 0 \\
Mixed-operation rows & 0 \\
Effective Levenshtein distance 3--5 & 0 \\
Another family closer & 0 \\
Gold tied nearest & 0 \\
Gold unique nearest & 0 \\
}

\newcommand{\FallbackRows}{%
Any unique nearest family & 161 & 14 & 125 & -111 & 98.0379\% \\
Only when A5 returns no result & 0 & 0 & 0 & +0 & 98.2055\% \\
Only queries of at most four characters & 13 & 0 & 9 & -9 & 98.1919\% \\
At most four characters and at most one bigram & 11 & 0 & 7 & -7 & 98.1949\% \\
}

\newcommand{\ManualFailureReviewRows}{%
Unique nearest, ranking & \code{soivinn} & \code{SOLVIN N} & \code{SOLVIN} & 20 & 1 & SOLVIN N (d=1, n=1) & SOLVIN N is uniquely nearest at distance 1, but the shorter SOLVIN family receives a much larger family score and pushes it to rank 20. \\
Nearest-distance tie & \code{conal} & \code{COBAL} & \code{CONIL} & 3 & 1 & COBAL; CONAZ; CONIL (d=1, n=3) & COBAL and CONIL are both one edit from CONAL, so raw distance cannot decide; A5 ranks COBAL third. \\
Another family closer & \code{rimunal} & \code{REMINYL} & \code{NIFUNAL} & 2 & 3 & NIFUNAL (d=2, n=1) & NIFUNAL is distance 2 while source target REMINYL is distance 3; A5's first result is the literal nearest family. \\
Another family closer & \code{blaxeedoil} & \code{BLACK SEED OIL} & \code{FLAX SEED OIL} & 2 & 3 & FLAX SEED OIL (d=2, n=1) & FLAX SEED OIL is distance 2 while BLACK SEED OIL is distance 3, making the generated phonetic corruption a real catalog near-collision. \\
Another family closer & \code{go5y} & \code{GOSAY} & \code{GORY} & 2 & 2 & GORY (d=1, n=1) & GORY is distance 1 while GOSAY is distance 2; the wrong first result follows the stronger visible spelling. \\
}

\newcommand{\UniqueScoreExampleRows}{%
\code{ligncainehydr0chloride} & \code{LIGNOCAINE HYDROCHLORIDE} & \code{LIDOCAINE HYDROCHLORIDE} & 2/3 & 1.8806/2.2115 & 2 & The uniquely nearest target is one edit closer, but the long-name neighbor leads the combined score by 0.3309. \\
\code{soivinn} & \code{SOLVIN N} & \code{SOLVIN} & 1/2 & 0.4585/1.5630 & 20 & Shorter-family evidence overwhelms the full SOLVIN N spelling by 1.1046 score units. \\
\code{clicynane} & \code{DICYNONE} & \code{CALCIONATE} & 3/4 & 1.0443/1.3082 & 9 & The structural CL-to-D target is one edit closer but remains at rank 9 after bounded corrections. \\
\code{ivyhuda} & \code{IVY BUTA} & \code{IVY HONTA} & 2/3 & 1.6709/1.7003 & 2 & The uniquely nearest target trails the displayed family by only 0.0294 score units. \\
}


\begin{document}
\begin{center}
{\LARGE\textbf{Clean Synthetic Core Medicine Search Benchmark}}\\[2pt]
{\large The 66,257-case primary test file, eleven retrieval systems, and failure evidence}\\[3pt]
{\small medicine-search-clean engineering analysis, 23 July 2026}
\end{center}
\vspace{0.25em}

\begin{center}
\fcolorbox{LineGray}{LightBlue}{%
\begin{minipage}{0.94\linewidth}
\textbf{Result in one paragraph.}
The main denominator is now \textbf{\DatasetCases} distinct, collision-free,
positive-retrieval query-target pairs from
\path{data/05_synthetic_clean_core/test_cases.csv}. Every system receives the
same case IDs and the verified target remains hidden until scoring. Algorithm 5
reaches \textbf{\AFiveHOne} Hit@1 and \textbf{\AFiveHTwenty} Hit@20. The table
results below are newly computed on this file; no score from the 595-row OCR or
115,000-row inclusive benchmark is reused.
\end{minipage}}
\end{center}

\section{Main Testing Contract}

\begin{table}[H]
\centering
\small
\begin{tabularx}{\linewidth}{@{}L{0.16\linewidth}L{0.29\linewidth}Y@{}}
\toprule
\textbf{Stage} & \textbf{Concrete example} & \textbf{What the benchmark does} \\
\midrule
Input & \code{0aflon} & Pass the corrupted commercial-name text to the search system. \\
Hidden answer & \code{DAFLON} & Keep the verified family outside retrieval; use it only after ranking. \\
Search output & ranked family names & Search the same 25,066-row runtime catalog with a limit of 20 families. \\
Scoring & Hit@1, Hit@5, Hit@10, Hit@20, MRR@20 & Locate the first acceptable family and calculate rank and safety metrics. \\
Primary vote & one case ID & Count each compact query-target pair once, regardless of source repetition. \\
\bottomrule
\end{tabularx}
\caption{Input, process, and output for every row. The expected family never enters candidate generation.}
\label{tab:contract}
\end{table}

\begin{table}[H]
\centering
\small
\begin{tabularx}{\linewidth}{@{}L{0.28\linewidth}rY@{}}
\toprule
\textbf{Audit item} & \textbf{Value} & \textbf{Meaning} \\
\midrule
\DatasetAuditRows
\bottomrule
\end{tabularx}
\caption{The clean-core CSV is the headline denominator. Counts reconcile from the 115,000-row source to retained source occurrences and then to unique primary pairs.}
\label{tab:audit}
\end{table}

\begin{table}[H]
\centering
\small
\begin{tabularx}{\linewidth}{@{}L{0.20\linewidth}L{0.21\linewidth}Y@{}}
\toprule
\textbf{Check} & \textbf{Observed result} & \textbf{Why it matters} \\
\midrule
Source identity & commit \code{47706ad}; SHA-256 \code{65f81b58dee1e7127386652383d2f6a8db1734e832dcbc73f14a9b831678886f} & The local CSV is byte-identical to the supplied clean-core file. \\
Case identity & \DatasetCases{} unique IDs; \DatasetCases{} unique compact queries & No row is counted twice. \\
Target coverage & \TargetFamilies{} unique verified families; zero unmapped targets & A retrieval miss cannot be caused by an absent expected family. \\
Expected behavior & all rows are \code{match} and \code{use\_in\_primary\_score=1} & Hit@k has one unambiguous positive-retrieval meaning. \\
Catalog status & all rows say \code{product\_scope\_unverified} & Product-scope validation remains a limitation even though every family key resolves. \\
Split & target-family hash, 80\% development and 20\% holdout & No target family occurs in both diagnostic splits. \\
\bottomrule
\end{tabularx}
\caption{Validation gates applied before any algorithm score was accepted.}
\label{tab:gates}
\end{table}

\section{Change From The 115,000-Row File}

The old file mixed positive typo retrieval with exact matches, no-match cases,
ambiguity behavior, route/form contradictions, and stress generators. The clean
core keeps only positive-retrieval rows with defensible text evidence, corrects
the effective mutation label, and collapses repeated compact query-target pairs.

\begin{table}[H]
\centering
\small
\begin{tabular}{@{}lrrr@{}}
\toprule
\textbf{Transformation} & \textbf{Before} & \textbf{After} & \textbf{Difference} \\
\midrule
Generated source rows & 115,000 & \RepresentedSourceRows & \RemovedSourceRows\ removed \\
Retained source occurrences & \RepresentedSourceRows & \DatasetCases & \CollapsedOccurrences\ duplicate occurrences collapsed \\
Primary scoring rows & 109,974 in the earlier fair view & \DatasetCases & New task-specific denominator \\
Exact real-name collision rows & 5,026 in the earlier fair audit & 0 & Excluded before this core was formed \\
\bottomrule
\end{tabular}
\caption{The clean core is not a filtered score view over 115,000 rows. It is a new unique-pair positive-retrieval benchmark.}
\label{tab:change-summary}
\end{table}

\begin{table}[H]
\centering
\small
\begin{tabular}{@{}lrrrrrr@{}}
\toprule
\textbf{System} & \textbf{Old fair H@1} & \textbf{Clean H@1} & \textbf{$\Delta$H@1} & \textbf{Old fair H@20} & \textbf{Clean H@20} & \textbf{$\Delta$H@20} \\
\midrule
\FairTransitionRows
\bottomrule
\end{tabular}
\caption{Algorithms 1--4 before and after changing the denominator. ``Old fair'' uses 109,974 collision-excluded generated rows; ``clean'' uses \DatasetCases{} unique positive-retrieval pairs. Delta measures dataset and task composition, not a code change.}
\label{tab:fair-transition}
\end{table}

\small
\begin{longtable}{@{}L{0.33\linewidth}rrL{0.40\linewidth}@{}}
\caption{Original category coverage after clean-core construction. A clean case can retain more than one provenance label, so membership counts are diagnostic rather than another denominator.}\label{tab:source-transition}\\
\toprule
\textbf{Original category} & \textbf{Source rows} & \textbf{Clean memberships} & \textbf{Role in the clean core} \\
\midrule
\endfirsthead
\toprule
\textbf{Original category} & \textbf{Source rows} & \textbf{Clean memberships} & \textbf{Role in the clean core} \\
\midrule
\endhead
\SourceTransitionRows
\bottomrule
\end{longtable}
\normalsize

\begin{table}[H]
\centering
\small
\begin{tabularx}{\linewidth}{@{}YrrL{0.35\linewidth}@{}}
\toprule
\textbf{Cleaning action} & \textbf{Rows} & \textbf{Share} & \textbf{Concrete interpretation} \\
\midrule
\CleaningNoteRows
\bottomrule
\end{tabularx}
\caption{Cleaning notes explain every relabelled or safety-retained row. Blank notes mean the original mutation label already matched the observed corruption.}
\label{tab:cleaning-notes}
\end{table}

\section{Dataset Composition}

\small
\begin{longtable}{@{}L{0.18\linewidth}L{0.26\linewidth}rrL{0.37\linewidth}@{}}
\caption{The dominant denominator contains one- and two-edit corruptions. Distance four and five remain as a small severe tail rather than dominating the score.}\label{tab:composition}\\
\toprule
\textbf{Dimension} & \textbf{Group} & \textbf{Cases} & \textbf{Share} & \textbf{Plain meaning} \\
\midrule
\endfirsthead
\toprule
\textbf{Dimension} & \textbf{Group} & \textbf{Cases} & \textbf{Share} & \textbf{Plain meaning} \\
\midrule
\endhead
\CompositionRows
\bottomrule
\end{longtable}
\normalsize

The source label \code{difficulty=EXTREME} contains
\GeneratorExtremeCases{} generator-design cases. It is not the normalized
distance cohort: the later \code{r>0.60} rule contains \ExtremeCases{} cases.

\small
\begin{longtable}{@{}L{0.34\linewidth}rrL{0.18\linewidth}L{0.18\linewidth}r@{}}
\caption{Clean mutation-category membership with one actual row-level example per category.}\label{tab:categories}\\
\toprule
\textbf{Clean category} & \textbf{Cases} & \textbf{Share} & \textbf{Example input} & \textbf{Expected family} & \textbf{Edits} \\
\midrule
\endfirsthead
\toprule
\textbf{Clean category} & \textbf{Cases} & \textbf{Share} & \textbf{Example input} & \textbf{Expected family} & \textbf{Edits} \\
\midrule
\endhead
\CategoryRows
\bottomrule
\end{longtable}
\normalsize

\small
\begin{longtable}{@{}L{0.39\linewidth}rrL{0.18\linewidth}L{0.22\linewidth}@{}}
\caption{Most frequent exact mutation labels. The complete machine-readable metric file contains all \ErrorTypeLabels{} labels.}\label{tab:error-types}\\
\toprule
\textbf{Clean error type} & \textbf{Cases} & \textbf{Share} & \textbf{Example input} & \textbf{Expected family} \\
\midrule
\endfirsthead
\toprule
\textbf{Clean error type} & \textbf{Cases} & \textbf{Share} & \textbf{Example input} & \textbf{Expected family} \\
\midrule
\endhead
\ErrorTypeRows
\bottomrule
\end{longtable}
\normalsize

\section{Methods And Metrics}

\small
\begin{longtable}{@{}rL{0.29\linewidth}L{0.61\linewidth}@{}}
\caption{The fixed comparison roster. Every row uses the same catalog, case IDs, top-20 cutoff, and scoring code.}\label{tab:systems}\\
\toprule
\textbf{Order} & \textbf{System} & \textbf{Evidence used} \\
\midrule
\endfirsthead
\toprule
\textbf{Order} & \textbf{System} & \textbf{Evidence used} \\
\midrule
\endhead
1 & Exact or prefix match & Exact compact family, then families that begin with the compact query. \\
2 & Exhaustive Levenshtein & Smallest insert, delete, or replace count over every family name. \\
3 & Jaro-Winkler & Character-order similarity with an explicit shared-prefix reward. \\
4 & Character 3-gram TF-IDF & Cosine similarity between weighted three-character pieces. \\
5 & RapidFuzz token ratio & Token-set and token-sort string similarity, keeping the larger score. \\
6 & Phonetic baseline & Medicine phonetic key, then Levenshtein ranking between sound keys. \\
7 & Algorithm 1 & Current application retrieval and safety ranking. \\
8 & Algorithm 2 & Fast indexed exact, edge, n-gram, skeleton, phonetic, and delete-key retrieval. \\
9 & Algorithm 3 & Weighted rank fusion of Algorithms 1 and 2. \\
10 & Algorithm 4 & Algorithm 2 plus family rescue, multi-signal scoring, bounded correction, and clarification. \\
11 & Algorithm 5 & Evidence-guided retrieval, mixed-error expansion, bounded reranking, and clarification. \\
\bottomrule
\end{longtable}
\normalsize

\subsection{Algorithm 5 Score Signals}

Algorithm 5 adds applicable signals rather than stopping at the first matching
rule. A weight is the signal's maximum score contribution, not a guarantee that
the candidate wins.

\small
\begin{longtable}{@{}L{0.19\linewidth}rL{0.30\linewidth}L{0.35\linewidth}@{}}
\caption{Every main signal in the Algorithm 5 rescue score.}\label{tab:lexical-signals}\\
\toprule
\textbf{Signal} & \textbf{Weight} & \textbf{Plain-English meaning} & \textbf{Example} \\
\midrule
\endfirsthead
\toprule
\textbf{Signal} & \textbf{Weight} & \textbf{Plain-English meaning} & \textbf{Example} \\
\midrule
\endhead
Exact family & 1.25 & Compact query exactly equals a catalog family. & \code{RIVOTRIL}$\rightarrow$\code{RIVOTRIL}. Exact evidence is strong; an exact different real drug is excluded from this clean core. \\
Ordinary edit similarity & 0.58 & Rewards fewer ordinary edits after accounting for string length. & Both \code{CEFAXIM} and \code{CEFIXIME} are one edit from \code{CEFAXIME}, so this signal ties. \\
Weighted edit similarity & 0.52 & Gives lower cost to known confusion pairs and vowel changes. Known pairs cost 0.45, vowel changes 0.70, and an ordinary replacement 1. & For \code{CEFAXIME}, \code{CEFIXIME} has weighted distance 0.70 while \code{CEFAXIM} has 1.00. \\
Same-position evidence & 0.24 & Measures query characters retained at the same indexes. & \code{CEFAXIME} and \code{CEFIXIME} match in seven of eight positions, producing 0.875. \\
Character 3-grams & 0.18 & Compares overlapping three-character pieces. & \code{RIVOTIL} and \code{RIVOTRIL} share \code{RIV}, \code{IVO}, and \code{VOT}. \\
Prefix evidence & 0.16 & Measures uninterrupted matching characters from the beginning. & \code{CLARANE} and \code{CLARA} share \code{CLARA}; this can favor a shorter wrong family. \\
Consonant skeleton & 0.16 & Compares a reduced key that preserves consonant structure. & \code{NOW} and \code{NEW} preserve \code{NW}; the key cannot decide which vowel was seen. \\
Length coverage & 0.16 & Rewards candidate length close to query length. & Eight-character \code{CEFIXIME} has full coverage; seven-character \code{CEFAXIM} has $7/8=0.875$. \\
Phonetic key & 0.14 & Compares sound-confusion groups such as C/K/Q, F/V, P/B, D/T, and S/Z. & \code{COUFSED} retains sound evidence for \code{COUGHSED}, but still needs lexical support. \\
Suffix evidence & 0.10 & Measures uninterrupted matching characters from the end. & \code{ANYOLAX} and \code{MYOLAX} share \code{YOLAX}. \\
Character 2-grams & 0.10 & Compares adjacent two-character pieces. & \code{MYMLAX} and \code{MYOLAX} retain several pairs despite one replacement. \\
Ordered subsequence & 0.10 & Rewards query letters that occur in candidate order despite gaps. & \code{RIVOTRI} is an ordered subsequence of \code{RIVOTRIL}. \\
\bottomrule
\end{longtable}
\normalsize

\subsection{Bonuses And Penalties}

\small
\begin{longtable}{@{}L{0.21\linewidth}rL{0.30\linewidth}L{0.37\linewidth}@{}}
\caption{Bonuses require supporting conditions. No single bonus is an automatic winner.}\label{tab:bonuses}\\
\toprule
\textbf{Condition} & \textbf{Change} & \textbf{Meaning} & \textbf{Specific word example} \\
\midrule
\endfirsthead
\toprule
\textbf{Condition} & \textbf{Change} & \textbf{Meaning} & \textbf{Specific word example} \\
\midrule
\endhead
Same first character & $+0.12$ & Useful evidence, too small to become a universal first-character rule. & \code{CEFAXIME} and \code{CEFIXIME} both start with C. \\
Confusable first characters & $+0.06$ & Retains a candidate when first letters form a known pair such as P/B or C/K. & \code{BANDOL}$\rightarrow$\code{PANDOL} receives the smaller B/P bonus. \\
Length difference at most 1 and weighted similarity at least 0.76 & $+0.18$ & Rewards a compact near-match only when spelling is already strong. & \code{CEFAXIME}$\rightarrow$\code{CEFIXIME} has equal length and weighted similarity 0.9125. \\
Prefix at least 0.55 and weighted similarity at least 0.66 & $+0.08$ & Rewards a strong beginning only when full spelling remains plausible. & \code{CLARANE}$\rightarrow$\code{CLARA} has prefix $5/7=0.714$. \\
Weighted similarity at least 0.84 and position at least 0.70 & $+0.22$ & Rewards agreement between independent spelling measurements. & \code{CEFAXIME}$\rightarrow$\code{CEFIXIME} passes both thresholds. \\
Short partial prefix with weak weighted evidence & $-0.18$ to $-0.42$ & Stops a short prefix from defeating a stronger complete name. & \code{ABIMOL}$\rightarrow$\code{ABIMOL EXTRA} can receive the maximum penalty. \\
Very short non-exact query with weak edit evidence & $-0.22$ & Reduces confidence from three- or four-character noise. & \code{ABC}$\rightarrow$\code{ABC IMMUNE} contains too little full-name evidence. \\
Weak spelling plus weak skeleton or phonetic evidence & $-0.20$ & Penalizes candidates supported by no strong lexical or sound channel. & \code{GHONT}$\rightarrow$\code{CALAMINE} has no strong support. \\
Weak prefix, weak 3-grams, and weak edit evidence & $-0.10$ & Penalizes little shared local structure. & \code{GHONT}$\rightarrow$\code{CALAMINE} also lacks a prefix and trigrams. \\
\bottomrule
\end{longtable}
\normalsize

\subsection{Bounded Corrections}

\small
\begin{longtable}{@{}L{0.15\linewidth}L{0.28\linewidth}L{0.22\linewidth}L{0.30\linewidth}@{}}
\caption{Except for controlled pure deletion, equal ordinary distance does not activate a correction.}\label{tab:corrections}\\
\toprule
\textbf{Correction} & \textbf{Required evidence} & \textbf{Why the bound exists} & \textbf{Specific word example} \\
\midrule
\endfirsthead
\toprule
\textbf{Correction} & \textbf{Required evidence} & \textbf{Why the bound exists} & \textbf{Specific word example} \\
\midrule
\endhead
Strictly closer candidate & Both retrieval layers find the alternative; distance at most 2 and at least one edit closer; score gap at most 0.25. & Agreement and the gap limit the change to a small inversion. & \code{MYCLAX}: supported \code{MYOLAX} can replace a distance-two first result. \\
Pure deletion candidate & Query length at least 5; ordered subsequence; one or two missing characters; no farther than current first; score gap at most 0.40. & Every visible character must stay in order. This is the only correction allowed at equal raw distance. & \code{RIVOTRI}$\rightarrow$\code{RIVOTRIL} is missing only the final L. \\
Single-token correction & Current first has several words; alternative has one; alternative is strictly closer within distance 2; gap at most 0.20. & A short family cannot win only because it is shorter. & \code{ABIMOL} may replace \code{ABIMOL EXTRA} when the full evidence is close. \\
Family-head correction & Catalog confirms a variant; query length at least 5; head within two edits and strictly closer; gap at most 1.40. & Catalog grouping blocks arbitrary prefixes from claiming family membership. & \code{LANTS} can recover the validated \code{LANTUS} variant family. \\
Strict full-name correction & Query length at least 5; rescue-only name within three edits and strictly closer; weighted distance no worse; gap at most 0.35; no protected current evidence. & The alternative must win on raw and weighted spelling without defeating independent protected evidence. & \code{MYOLANA} can move \code{MYOLAX} from rank 2 to rank 1. \\
\bottomrule
\end{longtable}
\normalsize

\begin{table}[H]
\centering
\small
\begin{tabularx}{\linewidth}{@{}L{0.18\linewidth}L{0.32\linewidth}Y@{}}
\toprule
\textbf{Metric} & \textbf{Calculation} & \textbf{Meaning} \\
\midrule
Hit@1 & target family occurs at rank 1 & The first displayed family is correct. \\
Hit@20 & target family occurs at ranks 1 through 20 & Retrieval found a usable candidate. \\
MRR@20 & $1/r$ for first target rank $r\leq20$, otherwise 0 & Earlier correct candidates receive more credit. \\
Failure rate & misses in group divided by rows in group & Finds the weakest group. \\
Failure share & misses in group divided by all misses for that system & Finds the group causing the most total damage. \\
Unsafe confident top-1 & wrong first family shown as confident without clarification & Safety failure; rank alone cannot determine it. \\
Latency & warm query wall time after index preparation & Implementation time on this machine, not asymptotic complexity. \\
\bottomrule
\end{tabularx}
\caption{Metric definitions used in every later table. Percentages never mix denominators.}
\label{tab:metrics}
\end{table}

\section{Overall Results}

\begin{table}[H]
\centering
\small
\begin{tabular}{@{}lrrrrrrrrrr@{}}
\toprule
\textbf{System} & \textbf{H@1} & \textbf{H@5} & \textbf{H@10} & \textbf{H@20} & \textbf{MRR@20} & \textbf{Build ms} & \textbf{Median ms} & \textbf{P95 ms} & \textbf{Mean candidates} & \textbf{Unsafe} \\
\midrule
\OverallRows
\bottomrule
\end{tabular}
\caption{Primary performance over all \DatasetCases{} clean unique pairs. Build time is one-time index preparation; median and P95 are warm query times.}
\label{tab:overall}
\end{table}

Accuracy rows use the same cases and scoring code. The refreshed Algorithm 5
latency was collected with eight parallel workers, while the retained baseline
latencies came from single-process runs. Those timing columns document each
artifact, but they are not a controlled speed comparison.

\begin{enumerate}
  \item \BestHOneSystem{} has the highest Hit@1 at \BestHOneScore;
        \BestHTwentySystem{} has the highest Hit@20 at \BestHTwentyScore.
  \item Relative to exhaustive Levenshtein, Algorithm 5 changes Hit@1 by
        \AFiveVsLevenHOne{} and Hit@20 by \AFiveVsLevenHTwenty.
\end{enumerate}

\begin{table}[H]
\centering
\small
\begin{tabular}{@{}lrrrr@{}}
\toprule
\textbf{System} & \textbf{Clarification rate} & \textbf{No-result rate} & \textbf{Unsafe confident top-1} & \textbf{Mean candidate count} \\
\midrule
\BehaviorRows
\bottomrule
\end{tabular}
\caption{Behavior uses response metadata, not rank alone. Mean candidate count is each algorithm's exposed considered or matched pool, not the common top-20 output size.}
\label{tab:behavior}
\end{table}

\begin{table}[H]
\centering
\small
\begin{tabular}{@{}lrrrrrr@{}}
\toprule
\textbf{System compared with A5} & \textbf{$\Delta$H@1} & \textbf{A5-only H@1} & \textbf{Other-only H@1} & \textbf{$p$ H@1} & \textbf{$\Delta$H@20} & \textbf{$p$ H@20} \\
\midrule
\PairedRows
\bottomrule
\end{tabular}
\caption{Paired case-level comparison. Delta is A5 minus the other system; exact McNemar $p$ tests whether discordant wins are balanced.}
\label{tab:paired}
\end{table}

\begin{table}[H]
\centering
\small
\begin{tabular}{@{}lrrrr@{}}
\toprule
\textbf{System} & \textbf{Dev. H@1} & \textbf{Holdout H@1} & \textbf{Dev. H@20} & \textbf{Holdout H@20} \\
\midrule
\SplitRows
\bottomrule
\end{tabular}
\caption{Target-family-disjoint diagnostic split. The all-case score remains primary; this table checks whether rankings transfer to held-out families.}
\label{tab:splits}
\end{table}

\section{Deep Audit Of The Remaining Rank-One Errors}

The strict pair-level score counts the source family as correct even when the
corrupted text is equally close or closer to another real catalog family. The
tables below preserve that score, then separate ranking defects, retrieval
defects, and text-level ambiguity.

\begin{table}[H]
\centering
\small
\begin{tabular}{@{}lrr@{}}
\toprule
\textbf{Mutually exclusive outcome} & \textbf{Cases} & \textbf{Share of 66,257} \\
\midrule
\ScoreReconciliationRows
\bottomrule
\end{tabular}
\caption{The arithmetic behind \AFiveHOne{} Hit@1. The \AFiveHOneFailures{}
first-rank misses contain \AFiveRankingFailures{} ranking failures, where the
expected family is still in the top 20, and \AFiveRetrievalFailures{} retrieval
failures. Only \AFiveNoResults{} queries return no candidates.}
\label{tab:failure-reconciliation}
\end{table}

\begin{table}[H]
\centering
\small
\begin{tabularx}{\linewidth}{@{}L{0.29\linewidth}rY@{}}
\toprule
\textbf{Independent check} & \textbf{Observed} & \textbf{Acceptance rule} \\
\midrule
\FailureValidityRows
\bottomrule
\end{tabularx}
\caption{The score was recomputed from row-level names and ranks. IDs, Hit@k
flags, Levenshtein values, unrestricted Damerau values, and exact-collision
checks reproduce exactly. The 88 stored-Damerau versus A5-OSA differences are
expected because those fields use different distance definitions; neither
field determines Hit@k. \FreshRerunCaption}
\label{tab:failure-validity}
\end{table}

The clean-core column \code{effective\_damerau} uses unrestricted
Damerau--Levenshtein distance. Algorithm 5's raw edit feature uses OSA, which
does not permit a character substring to be edited more than once. For example,
\code{altsren}$\rightarrow$\code{ALSOTRIN} has stored unrestricted-Damerau
distance 3 but A5 OSA distance 4. This definition difference occurs in 88 of
66,257 rows. It does not change the benchmark answer or the rank-based Hit@1
calculation.

\FreshRerunAuditText

\begin{table}[H]
\centering
\small
\begin{tabularx}{\linewidth}{@{}L{0.20\linewidth}rrrY@{}}
\toprule
\textbf{Denominator} & \textbf{Units} & \textbf{Hit@1} & \textbf{Hit@20} & \textbf{What receives one vote} \\
\midrule
\FamilyDenominatorRows
\bottomrule
\end{tabularx}
\caption{The headline is the pair micro-average. Family balancing lowers Hit@1 to \FamilyMacroHOne{} because small difficult families receive the same weight as families with many easy generated cases; it does not invalidate the pair-level score.}
\label{tab:family-denominators}
\end{table}

\begin{table}[H]
\centering
\small
\begin{tabular}{@{}lrr@{}}
\toprule
\textbf{Failure spread check} & \textbf{Value} & \textbf{Share} \\
\midrule
\FailureConcentrationRows
\bottomrule
\end{tabular}
\caption{Hit@1 failures occur across \FailureFamilies{} target families; the ten
largest family contributors create \TopTenFailureShare{} of all misses. The gap
is distributed, not produced by a short list of hard-coded names.}
\label{tab:failure-concentration}
\end{table}

\subsection{What The Corrupted Text Can Identify}

Algorithm 5 uses unweighted optimal-string-alignment distance, abbreviated OSA,
inside its raw edit feature. OSA charges one for insertion, deletion,
substitution, or one adjacent swap. For example, \code{firt} to \code{FRIT}
costs one OSA swap but two Levenshtein substitutions. The catalog scan compares
each query with all 17,476 compact family keys.

\begin{table}[H]
\centering
\small
\begin{tabular}{@{}lrrrrr@{}}
\toprule
\textbf{Catalog geometry under OSA} & \textbf{Cases} & \textbf{Share} & \textbf{A5 H@1} & \textbf{A5 H@20} & \textbf{H@1 misses} \\
\midrule
\CatalogGeometryRows
\bottomrule
\end{tabular}
\caption{A5 reaches \UniqueOSAHOne{} Hit@1 on the \UniqueOSACases{} cases where
the source target is the only nearest catalog family. Ties and closer real
families create \NonUniqueOSAFailures{} of the \AFiveHOneFailures{} strict Hit@1
misses.}
\label{tab:catalog-geometry}
\end{table}

\begin{table}[H]
\centering
\small
\begin{tabular}{@{}lrrrr@{}}
\toprule
\textbf{Status among A5 Hit@1 misses} & \textbf{Cases} & \textbf{Failure share} & \textbf{Ranking} & \textbf{Retrieval} \\
\midrule
\FailureRootRows
\bottomrule
\end{tabular}
\caption{Only \UniqueOSAFailures{} failures have a source target that is uniquely
nearest under A5's own raw distance. The other \NonUniqueOSAFailures{} are
catalog near-collisions or distance ties, although the source family remains the
strict benchmark answer.}
\label{tab:failure-roots}
\end{table}

\begin{table}[H]
\centering
\small
\begin{tabular}{@{}lrrrr@{}}
\toprule
\textbf{Status among \AFiveHOneFailures{} misses} & \textbf{Levenshtein $n$} & \textbf{Levenshtein share} & \textbf{OSA $n$} & \textbf{OSA share} \\
\midrule
\FailureDistanceDefinitionRows
\bottomrule
\end{tabular}
\caption{Levenshtein and OSA disagree on \FailureDistanceLabelDifferences{}
row-level failure labels because OSA counts an adjacent swap as one operation.
Both views are reported; OSA is primary for diagnosing A5 because it matches the
deployed raw-distance implementation.}
\label{tab:failure-distance-definitions}
\end{table}

\renewcommand{\arraystretch}{1.02}
\begin{table}[H]
\centering
\small
\begin{tabular}{@{}lrrrrrr@{}}
\toprule
\textbf{System} & \textbf{Unique H@1} & \textbf{Unique H@20} & \textbf{Tie H@1} & \textbf{Tie H@20} & \textbf{Closer H@1} & \textbf{Closer H@20} \\
\midrule
\SystemGeometryRows
\bottomrule
\end{tabular}
\caption{Performance changes sharply with catalog geometry. A5 is strongest on unique and tied targets, but a system cannot infer the generated source family from edit distance when another real family is literally closer.}
\label{tab:system-geometry}
\end{table}
\renewcommand{\arraystretch}{1.10}

\subsection{The \UniqueOSAFailures{} Unique-Nearest Failures}

These \UniqueOSAFailures{} rows are the clearest algorithm defects because no catalog family has
an equal or smaller OSA distance than the verified target.

\begin{table}[H]
\centering
\small
\begin{tabularx}{\linewidth}{@{}L{0.25\linewidth}rrY@{}}
\toprule
\textbf{Observed mechanism} & \textbf{Cases} & \textbf{Share of all A5 misses} & \textbf{Direct meaning} \\
\midrule
\UniqueFailureMechanismRows
\bottomrule
\end{tabularx}
\caption{The \UniqueOSAFailures{} unique-nearest failures divide into
\UniqueRankingFailures{} reranking defects and \UniqueRetrievalFailures{}
candidate-generation defects. The groups are exhaustive and mutually exclusive.}
\label{tab:unique-mechanisms}
\end{table}

\begin{table}[H]
\centering
\small
\begin{tabularx}{\linewidth}{@{}L{0.22\linewidth}rrrrY@{}}
\toprule
\textbf{Failure level} & \textbf{Cases} & \textbf{4+ bigrams} & \textbf{$\leq4$ chars} & \textbf{No result} & \textbf{Pattern} \\
\midrule
\UniqueFailurePatternRows
\bottomrule
\end{tabularx}
\caption{All \UniqueOSAFailures{} current unique-nearest failures are ranking
misses; no unique-nearest retrieval miss remains. Eight of these rows retain at
least four shared bigrams, so candidate order, not missing local evidence, is
the residual defect.}
\label{tab:unique-patterns}
\end{table}

\begin{table}[H]
\centering
\small
\begin{tabularx}{\linewidth}{@{}L{0.12\linewidth}L{0.15\linewidth}L{0.17\linewidth}rrrY@{}}
\toprule
\textbf{Input} & \textbf{Expected} & \textbf{A5 first} & \textbf{Raw $d$: gold/top} & \textbf{Score: gold/top} & \textbf{Rank} & \textbf{Manual diagnosis} \\
\midrule
\UniqueScoreExampleRows
\bottomrule
\end{tabularx}
\caption{Fresh A5 reruns expose the internal evidence on retained
unique-nearest ranking failures. These examples show why raw-distance advantage
can still lose to the combined score.}
\label{tab:unique-score-examples}
\end{table}

\subsection{Can Another Method Recover The Misses?}

\begin{table}[H]
\centering
\small
\begin{tabular}{@{}lrrrr@{}}
\toprule
\textbf{Other system} & \textbf{H@1 rescues} & \textbf{Share of A5 misses} & \textbf{H@20 retrievals} & \textbf{Share of A5 misses} \\
\midrule
\FailureRescueRows
\bottomrule
\end{tabular}
\caption{Each row asks how the other system handles exactly the cases A5 misses at rank 1. Counts overlap because several systems can recover the same case; exhaustive Levenshtein supplies the largest complementary H@1 set.}
\label{tab:failure-rescue}
\end{table}

\begin{table}[H]
\centering
\small
\begin{tabular}{@{}lr@{}}
\toprule
\textbf{Property of A5 retrieval misses that every other system also misses} & \textbf{Cases} \\
\midrule
\AllSystemMissRows
\bottomrule
\end{tabular}
\caption{Only \AllSystemMissCount{} current A5 retrieval misses also remain
outside every other system's top 20. A zero count confirms that Algorithm 5 now
retrieves every verified family in this clean core.}
\label{tab:all-system-misses}
\end{table}

\begin{table}[H]
\centering
\small
\begin{tabular}{@{}lrrrrr@{}}
\toprule
\textbf{Nearest-family override rule} & \textbf{Changed} & \textbf{Fixed} & \textbf{Broken} & \textbf{Net} & \textbf{Counterfactual H@1} \\
\midrule
\FallbackRows
\bottomrule
\end{tabular}
\caption{A global nearest-family override changes \FallbackAnyChanged{} rows,
fixes \FallbackAnyFixed{} misses, and breaks \FallbackAnyBroken{} correct
multi-signal decisions, a net change of \FallbackAnyNet{}. The no-result-only
fallback fixes \FallbackNoResultFixed{} and breaks \FallbackNoResultBroken{}
because the current algorithm has no empty results.}
\label{tab:fallback-counterfactual}
\end{table}

\subsection{Rows Inspected Manually}

\renewcommand{\arraystretch}{1.03}
\small
\begin{longtable}{@{}L{0.13\linewidth}L{0.08\linewidth}L{0.10\linewidth}L{0.11\linewidth}r r L{0.18\linewidth}L{0.24\linewidth}@{}}
\caption{\ManualReviewCount{} current failures inspected against the full catalog
and fresh A5 candidate evidence. Here $n$ is the number of catalog families at
the nearest OSA distance; $>20$ means absent from the returned top 20, not
literal rank 999.}\label{tab:manual-failure-review}\\
\toprule
\textbf{Class} & \textbf{Input} & \textbf{Expected} & \textbf{A5 first} & \textbf{Rank} & \textbf{Gold $d$} & \textbf{Nearest family, distance, count} & \textbf{What happened} \\
\midrule
\endfirsthead
\toprule
\textbf{Class} & \textbf{Input} & \textbf{Expected} & \textbf{A5 first} & \textbf{Rank} & \textbf{Gold $d$} & \textbf{Nearest family, distance, count} & \textbf{What happened} \\
\midrule
\endhead
\ManualFailureReviewRows
\bottomrule
\end{longtable}
\normalsize
\renewcommand{\arraystretch}{1.10}

\section{Which Errors Cause Failure}

\begin{table}[H]
\centering
\small
\begin{tabular}{@{}lrrrrrr@{}}
\toprule
\textbf{System} & \textbf{Overall H@1} & \textbf{1 edit} & \textbf{2 edits} & \textbf{3 edits} & \textbf{4 edits} & \textbf{5 edits} \\
\midrule
\DistanceHOneRows
\bottomrule
\end{tabular}
\caption{Hit@1 by effective Levenshtein count. One edit means one insertion, deletion, or replacement; adjacent swaps still cost two here.}
\label{tab:distance-h1}
\end{table}

\begin{table}[H]
\centering
\small
\begin{tabular}{@{}lrrrrrr@{}}
\toprule
\textbf{System} & \textbf{Overall H@20} & \textbf{1 edit} & \textbf{2 edits} & \textbf{3 edits} & \textbf{4 edits} & \textbf{5 edits} \\
\midrule
\DistanceHTwentyRows
\bottomrule
\end{tabular}
\caption{Hit@20 by effective Levenshtein count. The distance-four and distance-five columns expose candidate-generation limits that overall accuracy can hide.}
\label{tab:distance-h20}
\end{table}

\begin{table}[H]
\centering
\small
\begin{tabular}{@{}lrrrrrrrrr@{}}
\toprule
\textbf{System} & \textbf{Mixed} & \textbf{Visual/OCR} & \textbf{Phonetic} & \textbf{Vowel} & \textbf{Deletion} & \textbf{Swap} & \textbf{Insertion} & \textbf{Keyboard} & \textbf{Other} \\
\midrule
\OperationRows
\bottomrule
\end{tabular}
\caption{Hit@1 by rule-derived mutation mechanism. Mixed means that category
and error labels contain more than one mechanism, such as deletion plus
phonetic substitution. Hit@20 is 100\% for Algorithm 5 in every mechanism.}
\label{tab:operations}
\end{table}

\begin{center}
\fcolorbox{LineGray}{LightYellow}{%
\begin{minipage}{0.94\linewidth}
\textbf{What is a bigram?}
A character bigram is two adjacent compact characters. \code{DAFLON} contains
\code{DA}, \code{AF}, \code{FL}, \code{LO}, and \code{ON}. Query
\code{0aflon} shares \code{AF}, \code{FL}, \code{LO}, and \code{ON} with
\code{DAFLON}. A zero-bigram row keeps no adjacent two-character anchor.
\end{minipage}}
\end{center}

\begin{table}[H]
\centering
\small
\begin{tabular}{@{}lrrrrr@{}}
\toprule
\textbf{System} & \textbf{Overall H@1} & \textbf{0 shared} & \textbf{1 shared} & \textbf{2--3 shared} & \textbf{4+ shared} \\
\midrule
\BigramRows
\bottomrule
\end{tabular}
\caption{Hit@1 by shared adjacent-character evidence. The table measures how
much local text remains for first-rank ordering; Algorithm 5 still retrieves all
groups within the top 20.}
\label{tab:bigrams}
\end{table}

Algorithm 5 reaches \AFiveZeroBigramHOne{} Hit@1 when the query and target share
no adjacent pair. This is the strictest first-rank test with almost no local
lexical anchor; Hit@20 remains 100\%.

For compact query $q_c$ and compact target $y_c$, normalized distance is
\[
  r=\frac{L(q_c,y_c)}{\max(|y_c|,1)},
  \qquad \text{extreme when }r>0.60,
\]
where $L$ is Levenshtein distance. The threshold means more than 60 edits per
100 target characters, not 60\% of the dataset. For example,
\code{2ix}$\rightarrow$\code{ZEX} has two edits over a three-character target,
so $r=2/3=0.667$ and the row is extreme.

\begin{table}[H]
\centering
\small
\begin{tabular}{@{}lrrrrrrrrr@{}}
\toprule
\textbf{System} & \textbf{Std. $n$} & \textbf{Std. H@1} & \textbf{Std. H@20} & \textbf{High $n$} & \textbf{High H@1} & \textbf{High H@20} & \textbf{Extreme $n$} & \textbf{Extreme H@1} & \textbf{Extreme H@20} \\
\midrule
\CohortRows
\bottomrule
\end{tabular}
\caption{Recovery by normalized-distance cohort. Standard has $r\leq0.40$, high has $0.40<r\leq0.60$, and extreme has $r>0.60$.}
\label{tab:cohorts}
\end{table}

The extreme cohort contains \ExtremeCases{} cases. Algorithm 5 reaches
\AFiveExtremeHOne{} Hit@1 and \AFiveExtremeHTwenty{} Hit@20 on those rows.

\small
\begin{longtable}{@{}L{0.20\linewidth}L{0.27\linewidth}rrrr@{}}
\caption{Distribution and Algorithm 5 recovery inside the extreme cohort. Rows under different dimensions describe the same extreme cases and should not be summed together.}\label{tab:extreme-distribution}\\
\toprule
\textbf{Dimension} & \textbf{Group} & \textbf{Cases} & \textbf{Share} & \textbf{A5 H@1} & \textbf{A5 H@20} \\
\midrule
\endfirsthead
\toprule
\textbf{Dimension} & \textbf{Group} & \textbf{Cases} & \textbf{Share} & \textbf{A5 H@1} & \textbf{A5 H@20} \\
\midrule
\endhead
\ExtremeDistributionRows
\bottomrule
\end{longtable}
\normalsize

\begin{table}[H]
\centering
\small
\begin{tabularx}{\linewidth}{@{}L{0.13\linewidth}L{0.15\linewidth}rrL{0.16\linewidth}rY@{}}
\toprule
\textbf{Input} & \textbf{Expected} & \textbf{$r$} & \textbf{Shared bigrams} & \textbf{A5 first} & \textbf{Expected rank} & \textbf{Category} \\
\midrule
\ExtremeExampleRows
\bottomrule
\end{tabularx}
\caption{Concrete extreme rows show what information survives, what A5 returns, and whether the verified family remains inside the top 20.}
\label{tab:extreme-examples}
\end{table}

\small
\begin{longtable}{@{}L{0.30\linewidth}rrrrrrr@{}}
\caption{Hit@20 by clean category for all six classical baselines. This exposes which mutation mechanisms suit each single-signal method.}\label{tab:classical-category}\\
\toprule
\textbf{Clean category} & \textbf{$n$} & \textbf{Exact/prefix} & \textbf{Leven.} & \textbf{Jaro-W.} & \textbf{3-gram} & \textbf{Token} & \textbf{Phonetic} \\
\midrule
\endfirsthead
\toprule
\textbf{Clean category} & \textbf{$n$} & \textbf{Exact/prefix} & \textbf{Leven.} & \textbf{Jaro-W.} & \textbf{3-gram} & \textbf{Token} & \textbf{Phonetic} \\
\midrule
\endhead
\ClassicalCategoryRows
\bottomrule
\end{longtable}
\normalsize

\small
\begin{longtable}{@{}L{0.25\linewidth}rrrrrrrrrrr@{}}
\caption{Hit@1 and Hit@20 by clean category for Algorithms 1--5. Each row reports the same category memberships and cases for all systems.}\label{tab:category-h1}\\
\toprule
\textbf{Clean category} & \textbf{$n$} & \textbf{A1 H@1} & \textbf{A2 H@1} & \textbf{A3 H@1} & \textbf{A4 H@1} & \textbf{A5 H@1} & \textbf{A1 H@20} & \textbf{A2 H@20} & \textbf{A3 H@20} & \textbf{A4 H@20} & \textbf{A5 H@20} \\
\midrule
\endfirsthead
\toprule
\textbf{Clean category} & \textbf{$n$} & \textbf{A1 H@1} & \textbf{A2 H@1} & \textbf{A3 H@1} & \textbf{A4 H@1} & \textbf{A5 H@1} & \textbf{A1 H@20} & \textbf{A2 H@20} & \textbf{A3 H@20} & \textbf{A4 H@20} & \textbf{A5 H@20} \\
\midrule
\endhead
\CategoryAlgorithmRows
\bottomrule
\end{longtable}
\normalsize

\renewcommand{\arraystretch}{1.00}
\small
\begin{longtable}{@{}L{0.34\linewidth}rrrrrr@{}}
\caption{Algorithm 5 Hit@1 failure impact by clean category. Membership
categories can overlap on rare rows, so shares are not summed.
\AFiveWorstCategory{} creates the largest absolute first-rank miss count,
\AFiveWorstCategoryMisses{} cases; this need not be the highest within-category
failure rate.}\label{tab:a5-impact}\\
\toprule
\textbf{Clean category} & \textbf{Cases} & \textbf{H@1} & \textbf{H@20} & \textbf{Misses} & \textbf{Failure rate} & \textbf{Failure share} \\
\midrule
\endfirsthead
\toprule
\textbf{Clean category} & \textbf{Cases} & \textbf{H@1} & \textbf{H@20} & \textbf{Misses} & \textbf{Failure rate} & \textbf{Failure share} \\
\midrule
\endhead
\AFiveImpactRows
\bottomrule
\end{longtable}
\normalsize
\renewcommand{\arraystretch}{1.10}

\section{Danger And Equal-Distance Cases}

\small
\begin{longtable}{@{}lrrrr@{}}
\caption{Performance on \DangerCases{} gold-closer dangerous pairs. The query is closer to the expected family than the named competitor, so these remain scoreable but need cautious presentation.}\label{tab:danger}\\
\toprule
\textbf{System} & \textbf{Danger cases} & \textbf{H@1} & \textbf{H@20} & \textbf{Unsafe confident top-1} \\
\midrule
\endfirsthead
\toprule
\textbf{System} & \textbf{Danger cases} & \textbf{H@1} & \textbf{H@20} & \textbf{Unsafe confident top-1} \\
\midrule
\endhead
\DangerRows
\bottomrule
\end{longtable}
\normalsize

\begin{table}[H]
\centering
\small
\begin{tabularx}{\linewidth}{@{}L{0.15\linewidth}L{0.15\linewidth}L{0.15\linewidth}rrY@{}}
\toprule
\textbf{Input} & \textbf{Expected} & \textbf{Recorded competitor} & \textbf{Gold edits} & \textbf{Competitor edits} & \textbf{Why it remains} \\
\midrule
\DangerExampleRows
\bottomrule
\end{tabularx}
\caption{Concrete dangerous rows. Unlike an exact real-drug collision, the expected family has strictly smaller edit distance than the alternative.}
\label{tab:danger-examples}
\end{table}

\begin{table}[H]
\centering
\small
\begin{tabular}{@{}lrrrrr@{}}
\toprule
\textbf{System} & \textbf{Wrong top-1 with result} & \textbf{Equal-distance wrong top-1} & \textbf{Share} & \textbf{Rank-only ties} & \textbf{Retrieval ties} \\
\midrule
\EqualDistanceRows
\bottomrule
\end{tabular}
\caption{A tie means the wrong first family and verified family have the same raw Levenshtein distance from the query. Equal distance identifies ambiguity, not the correct winner.}
\label{tab:equal-distance}
\end{table}

\begin{table}[H]
\centering
\small
\begin{tabularx}{\linewidth}{@{}L{0.13\linewidth}L{0.16\linewidth}L{0.16\linewidth}rrL{0.23\linewidth}Y@{}}
\toprule
\textbf{Input} & \textbf{Expected} & \textbf{A5 first} & \textbf{Tie distance} & \textbf{Expected rank} & \textbf{Category} & \textbf{Failure level} \\
\midrule
\EqualDistanceExampleRows
\bottomrule
\end{tabularx}
\caption{Actual Algorithm 5 equal-distance failures from the clean core. Ranking failures already contain the target in the top 20; retrieval failures do not.}
\label{tab:tie-examples}
\end{table}

\begin{table}[H]
\centering
\small
\begin{tabularx}{\linewidth}{@{}cL{0.24\linewidth}Y@{}}
\toprule
\textbf{Order} & \textbf{Equal-distance handling} & \textbf{Operational rule} \\
\midrule
1 & Keep the complete A5 score order & The score combines raw and weighted edits, position, edge, n-gram, skeleton, phonetic, length, family, and retriever agreement. \\
2 & Permit only bounded corrections & A candidate moves first only when a documented correction has enough independent evidence and a limited score gap. \\
3 & Mark the response ambiguous & Equal raw distance does not justify hidden certainty; show alternatives. \\
4 & Ask for visible evidence & Use strength, form, route, or unreadable-position information to select the family or variant. \\
\bottomrule
\end{tabularx}
\caption{The ranking policy is evidence aggregation plus clarification, not first-letter or alphabetic priority.}
\label{tab:tie-policy}
\end{table}

\section{Concrete Inputs And Outputs}

\begin{table}[H]
\centering
\small
\begin{tabularx}{\linewidth}{@{}L{0.22\linewidth}L{0.13\linewidth}L{0.16\linewidth}L{0.16\linewidth}rY@{}}
\toprule
\textbf{Comparison} & \textbf{Input} & \textbf{Expected} & \textbf{Other top-1} & \textbf{Other rank} & \textbf{Why A5 adds value} \\
\midrule
\AFiveGainRows
\bottomrule
\end{tabularx}
\caption{One real A5-only Hit@1 case against each comparison system when such a case exists. These rows expose complementary retrieval rather than average accuracy alone.}
\label{tab:a5-gains}
\end{table}

\small
\begin{longtable}{@{}L{0.20\linewidth}L{0.08\linewidth}L{0.11\linewidth}L{0.12\linewidth}L{0.12\linewidth}rrL{0.17\linewidth}@{}}
\caption{Up to one ranking failure and one retrieval failure per system, selected deterministically when that failure type exists in the canonical row-level output.}\label{tab:failures}\\
\toprule
\textbf{System} & \textbf{Failure} & \textbf{Input} & \textbf{Expected} & \textbf{Top-1} & \textbf{Rank} & \textbf{Edits} & \textbf{Category} \\
\midrule
\endfirsthead
\toprule
\textbf{System} & \textbf{Failure} & \textbf{Input} & \textbf{Expected} & \textbf{Top-1} & \textbf{Rank} & \textbf{Edits} & \textbf{Category} \\
\midrule
\endhead
\FailureExampleRows
\bottomrule
\end{longtable}
\normalsize

\section{Engineering Decision}

\begin{table}[H]
\centering
\small
\begin{tabularx}{\linewidth}{@{}cL{0.24\linewidth}L{0.31\linewidth}Y@{}}
\toprule
\textbf{Order} & \textbf{Action} & \textbf{Evidence in this report} & \textbf{Acceptance condition} \\
\midrule
P0 & Adopt the clean core as the primary synthetic retrieval file & Every one of \DatasetCases{} rows is unique, catalog-resolvable, and scoreable. & Future systems run all case IDs with no silent skips. \\
P1 & Improve the highest-impact A5 failure groups & Table~\ref{tab:a5-impact} separates rate from total miss contribution. & Raise holdout H@20 without increasing unsafe confidence. \\
P2 & Preserve ambiguity for equal raw distance & Table~\ref{tab:equal-distance} shows ties remain common among wrong first results. & Alternatives remain visible; no first-character hard code. \\
P3 & Validate product scope and dangerous alternatives & Every row is currently marked product-scope unverified. & A domain reviewer confirms target and alternative identities against the deployment catalog. \\
P4 & Keep OCR and human-study claims separate & This file contains generated text queries, not prescription images or pharmacist decisions. & OCR and user-study results retain their own denominators. \\
\bottomrule
\end{tabularx}
\caption{Ordered engineering work derived from the measured failure distribution.}
\label{tab:actions}
\end{table}

\end{document}
